% Lesson 7 — Reading: Texts on recreational activities — Anglais Tle A — Cours signé OnBuch+
% Programme officiel MINESEC. Compiler avec : tectonic EN07-reading-texts-on-recreational-activities.tex
\newcommand{\DOCMATIERE}{Anglais}
\newcommand{\DOCNIVEAU}{Tle A}
\newcommand{\DOCMODULE}{Chapter 1 : Family and Social Life}
\newcommand{\DOCLECON}{EN07}
\newcommand{\DOCTITRE}{Lesson 7 — Reading: Texts on recreational activities}
% OnBuch+ — préambule « cours signés » ENRICHI (compilé avec tectonic / XeLaTeX)
% Système visuel « Pop » d'OnBuch+ : fond crème (jamais blanc), bordures encre
% épaisses, ombres « dures » décalées, titres Archivo Black en capitales.
% Version MATHÉMATIQUES : graphes (pgfplots/tikz), boîtes pédagogiques
% (définition, propriété, méthode, attention, le savais-tu, exercice résolu,
% exercice, TP, bilan), page de garde et signature OnBuch+ ancrée en bas.
%
% Macros attendues dans main.tex AVANT \input{preamble} :
%   \DOCMATIERE  \DOCNIVEAU  \DOCMODULE  \DOCLECON (numéro)  \DOCTITRE
\documentclass[11pt]{article}

\usepackage{fontspec}
\usepackage[english,french]{babel} % français = langue principale ; l'anglais via \eng{..} / textebox
\usepackage[a4paper,margin=2cm,top=3.4cm,bottom=2.8cm,headheight=1.4cm,footskip=1.3cm]{geometry}
\usepackage{amsmath,amssymb,amsfonts}
\usepackage{mathtools} % \xmapsto, \coloneqq... souvent utilisés par les modèles
\usepackage{cancel} % simplifications barrées (bilans de forces)
% \abs{z} (module, valeur absolue) et \norm{u} (norme), utilisés par les modèles
\providecommand{\abs}{}\let\abs\relax\DeclarePairedDelimiter{\abs}{\lvert}{\rvert}
\providecommand{\norm}{}\let\norm\relax\DeclarePairedDelimiter{\norm}{\lVert}{\rVert}
\usepackage[table]{xcolor} % [table] : fournit aussi \rowcolors (alternance de lignes)
\usepackage{pagecolor}
\usepackage{tikz}
\usetikzlibrary{arrows.meta,positioning,calc,shapes.geometric,shapes.misc,shapes.arrows,decorations.pathmorphing,decorations.pathreplacing,decorations.markings,patterns,shadows,mindmap,trees,fit,backgrounds,matrix,intersections,angles,quotes}
\usepackage{pgfplots}
\usepackage[version=4]{mhchem} % équations nucléaires \ce{^{235}_{92}U}
\usepackage[european,siunitx]{circuitikz} % schémas électriques
\pgfplotsset{compat=1.18}
\usepgfplotslibrary{fillbetween,groupplots} % aires entre courbes (calcul intégral)
\usepackage{enumitem}
\usepackage{fancyhdr}
% [most] charge toutes les bibliothèques tcolorbox (listings, minted, raster,
% documentation...) et gonfle inutilement la mémoire TeX sur un document long
% (~300 boîtes) : on ne charge que ce qui est réellement utilisé.
\usepackage[skins,breakable]{tcolorbox}
\usepackage{graphicx}
\usepackage{colortbl}
\usepackage{booktabs}
\usepackage{tabularx}
\usepackage{diagbox} % case d'en-tête diagonale des tableaux à double entrée
% Colonnes extensibles centrée / gauche / droite (C, L, R), souvent utilisées par les modèles
\newcolumntype{C}{>{\centering\arraybackslash}X}
\newcolumntype{L}{>{\raggedright\arraybackslash}X}
\newcolumntype{R}{>{\raggedleft\arraybackslash}X}
\usepackage{array}
\usepackage{multicol}
\usepackage{multirow}
\usepackage{siunitx}
% parse-numbers=false : accepte aussi \SI{2,5 \times 10^{-3}}{...}
\sisetup{locale=FR,output-decimal-marker={,},group-minimum-digits=4,parse-numbers=false}
\DeclareSIUnit\ohm{\ensuremath{\symup{\Omega}}} % U+2126 absent de la police
\DeclareSIUnit\division{div} % réglages d'oscilloscope (ms/div, V/div)
\DeclareSIUnit\tr{tr} % tours (vitesse de rotation, ex. \SI{1500}{\tr\per\minute})
\DeclareSIUnit\molar{M} % concentration molaire (ex. \SI{3}{\molar}, \SI{1}{\micro\molar})
\DeclareSIUnit\an{an} % année (le modèle confond parfois avec \year, un compteur TeX) — ex. \SI{5}{\kilo\meter\per\an}
\DeclareSIUnit\FCFA{FCFA} % franc CFA (ex. \SI{500}{\FCFA\per\kilo\gram})
\DeclareSIUnit\degreeBrix{\textdegree Brix} % teneur en sucre d'un sirop/jus (réfractométrie)
\DeclareSIUnit\month{mois} % le modèle utilise parfois \month, un compteur TeX, comme unité de durée
\usepackage{qrcode}
% \degree : commande fréquemment utilisée par les modèles pour un angle ou une
% température en dehors de \SI{}{} (non fournie par les paquets chargés).
\providecommand{\degree}{^{\circ}}
% \qed (fin de démonstration), utilisé par les modèles sans amsthm
\providecommand{\qed}{\hfill$\square$}
% \barre : double barre des valeurs interdites dans un tableau de variations (tkz-tab)
\providecommand{\barre}{\ensuremath{\|}}
% environnement proof (amsthm non chargé) utilisé par les modèles
\ifdefined\proof\else\newenvironment{proof}[1][Démonstration]{\par\noindent\textit{#1.}\ }{\qed\par}\fi
\usepackage{lastpage}
\usepackage{hyperref}
\hypersetup{hidelinks}

% ============================================================
% Palette « Pop » OnBuch+ — mêmes hex que la classe P (plus_ui.dart)
% ============================================================
\definecolor{poporange}{HTML}{F59321}
\definecolor{popdark}{HTML}{E07A0C}
\definecolor{popcream}{HTML}{FFF6E8}
\definecolor{popink}{HTML}{1C1714}
\definecolor{poppurple}{HTML}{7A5AE0}
\definecolor{popgreen}{HTML}{2E9E6A}
\definecolor{poppink}{HTML}{E0457B}
\definecolor{popblue}{HTML}{2F7DE1}
\definecolor{popgold}{HTML}{C98A12}
\definecolor{popmuted}{HTML}{8A7F76}
\definecolor{popline}{HTML}{EADFD2}
\definecolor{popgray}{HTML}{8A7F76}
\colorlet{popgrey}{popgray}
% Alias tolérants (le modèle nomme parfois les couleurs différemment)
\colorlet{popwhite}{white}
\colorlet{popblack}{popink}
\colorlet{popred}{poppink}
\colorlet{popyellow}{popgold}
% Teintes claires (fonds de tableaux/encadrés) — le modèle nomme parfois
% les couleurs avec un suffixe L (ex. popblueL) sans les avoir définies.
\colorlet{poporangeL}{poporange!20}
\colorlet{popdarkL}{popdark!20}
\colorlet{poppurpleL}{poppurple!20}
\colorlet{popgreenL}{popgreen!20}
\colorlet{poppinkL}{poppink!20}
\colorlet{popblueL}{popblue!20}
\colorlet{popgoldL}{popgold!20}
\colorlet{popredL}{popred!20}
\colorlet{popgrayL}{popgray!20}
% Teintes claires pour fonds de tableaux / zones
\colorlet{popblueL}{popblue!10!white}
\colorlet{poporangeL}{poporange!14!white}
\colorlet{popgreenL}{popgreen!12!white}
\colorlet{poppurpleL}{poppurple!12!white}
\colorlet{poppinkL}{poppink!12!white}
\colorlet{popgoldL}{popgold!14!white}

\pagecolor{popcream}

% ============================================================
% Typographie
% ============================================================
\defaultfontfeatures{Ligatures=TeX}
\setmainfont{PlusJakartaSans-Medium.ttf}[Path=../../../fonts/,BoldFont=PlusJakartaSans-Bold.ttf]
% Maths en sans-serif (Fira Math) avec chiffres et lettres droites de Plus
% Jakarta, pour que les formules se fondent dans le texte.
\usepackage[math-style=ISO,bold-style=ISO]{unicode-math}
\setmathfont{FiraMath-Regular.otf}[Path=../../../fonts/]
\setmathfont{PlusJakartaSans-Medium.ttf}[Path=../../../fonts/,range={up/num,up/{Latin,latin},\mathup}]
\usepackage{icomma} % virgule décimale sans espace en mode math
% Caractères mathématiques Unicode absents des polices de texte (Plus Jakarta,
% Archivo Black) : rendus par leur équivalent mathématique, en texte comme en
% formule — sinon ils s'affichent en carré vide (ex. « Arithmétique dans ℤ »).
\usepackage{newunicodechar}
\newunicodechar{ℝ}{\ensuremath{\mathbb{R}}}
\newunicodechar{ℤ}{\ensuremath{\mathbb{Z}}}
\newunicodechar{ℂ}{\ensuremath{\mathbb{C}}}
\newunicodechar{ℕ}{\ensuremath{\mathbb{N}}}
\newunicodechar{ℚ}{\ensuremath{\mathbb{Q}}}
\newunicodechar{𝔻}{\ensuremath{\mathbb{D}}}
\newunicodechar{α}{\ensuremath{\alpha}}
\newunicodechar{β}{\ensuremath{\beta}}
\newunicodechar{γ}{\ensuremath{\gamma}}
\newunicodechar{δ}{\ensuremath{\delta}}
\newunicodechar{ε}{\ensuremath{\varepsilon}}
\newunicodechar{θ}{\ensuremath{\theta}}
\newunicodechar{λ}{\ensuremath{\lambda}}
\newunicodechar{μ}{\ensuremath{\mu}}
\newunicodechar{σ}{\ensuremath{\sigma}}
\newunicodechar{τ}{\ensuremath{\tau}}
\newunicodechar{φ}{\ensuremath{\varphi}}
\newunicodechar{ψ}{\ensuremath{\psi}}
\newunicodechar{ω}{\ensuremath{\omega}}
\newunicodechar{Σ}{\ensuremath{\Sigma}}
% majuscules produites par \MakeUppercase dans les titres (α-aminés -> Α)
\newunicodechar{Α}{\ensuremath{\alpha}}
\newunicodechar{Β}{\ensuremath{\beta}}
\newunicodechar{Γ}{\ensuremath{\Gamma}}
\newunicodechar{Δ}{\ensuremath{\Delta}}
\newunicodechar{Ε}{\ensuremath{\varepsilon}}
\newunicodechar{Θ}{\ensuremath{\Theta}}
\newunicodechar{Λ}{\ensuremath{\Lambda}}
\newunicodechar{Μ}{\ensuremath{\mu}}
\newunicodechar{Τ}{\ensuremath{\tau}}
\newunicodechar{Φ}{\ensuremath{\Phi}}
\newunicodechar{Ψ}{\ensuremath{\Psi}}
\newunicodechar{Ω}{\ensuremath{\Omega}}
\newunicodechar{↦}{\ensuremath{\mapsto}}
\newunicodechar{∈}{\ensuremath{\in}}
\newunicodechar{∉}{\ensuremath{\notin}}
\newunicodechar{∧}{\ensuremath{\wedge}}
\newunicodechar{⇔}{\ensuremath{\Leftrightarrow}}
\newunicodechar{⟺}{\ensuremath{\Longleftrightarrow}}
\newunicodechar{⇒}{\ensuremath{\Rightarrow}}
\newunicodechar{⟹}{\ensuremath{\Longrightarrow}}
\newunicodechar{∘}{\ensuremath{\circ}}
\newunicodechar{∪}{\ensuremath{\cup}}
\newunicodechar{∩}{\ensuremath{\cap}}
\newunicodechar{≡}{\ensuremath{\equiv}}
\newunicodechar{⊂}{\ensuremath{\subset}}
\newunicodechar{⊥}{\ensuremath{\perp}}
\newunicodechar{∃}{\ensuremath{\exists}}
\newunicodechar{∀}{\ensuremath{\forall}}
\newunicodechar{∛}{\ensuremath{\sqrt[3]{\,}}}
\newunicodechar{∜}{\ensuremath{\sqrt[4]{\,}}}
\newunicodechar{⃗}{\ensuremath{{}^{\to}}}
\newunicodechar{ⁿ}{\ifmmode^{n}\else\textsuperscript{\ensuremath{n}}\fi}
\newunicodechar{ˣ}{\ifmmode^{x}\else\textsuperscript{\ensuremath{x}}\fi}
\newunicodechar{ᵇ}{\ifmmode^{b}\else\textsuperscript{\ensuremath{b}}\fi}
\newunicodechar{ʳ}{\ifmmode^{r}\else\textsuperscript{\ensuremath{r}}\fi}
\newunicodechar{ᵘ}{\ifmmode^{u}\else\textsuperscript{\ensuremath{u}}\fi}
\newunicodechar{ʸ}{\ifmmode^{y}\else\textsuperscript{\ensuremath{y}}\fi}
\newunicodechar{ᵃ}{\ifmmode^{a}\else\textsuperscript{\ensuremath{a}}\fi}
\newunicodechar{ᵉ}{\ifmmode^{e}\else\textsuperscript{\ensuremath{e}}\fi}
\newunicodechar{ᵏ}{\ifmmode^{k}\else\textsuperscript{\ensuremath{k}}\fi}
\newunicodechar{ᵖ}{\ifmmode^{p}\else\textsuperscript{\ensuremath{p}}\fi}
\newunicodechar{ᴺ}{\ifmmode^{N}\else\textsuperscript{\ensuremath{N}}\fi}
\newunicodechar{ᶜ}{\ifmmode^{c}\else\textsuperscript{\ensuremath{c}}\fi}
\newunicodechar{ʰ}{\ifmmode^{h}\else\textsuperscript{\ensuremath{h}}\fi}
\newunicodechar{ᵠ}{\ifmmode^{q}\else\textsuperscript{\ensuremath{q}}\fi}
\newunicodechar{ₙ}{\ifmmode_{n}\else\textsubscript{\ensuremath{n}}\fi}
\newunicodechar{ₐ}{\ifmmode_{a}\else\textsubscript{\ensuremath{a}}\fi}
\newunicodechar{ᵢ}{\ifmmode_{i}\else\textsubscript{\ensuremath{i}}\fi}
\newunicodechar{ᵣ}{\ifmmode_{r}\else\textsubscript{\ensuremath{r}}\fi}
\newunicodechar{ₖ}{\ifmmode_{k}\else\textsubscript{\ensuremath{k}}\fi}
\newunicodechar{ᵦ}{\ifmmode_{\beta}\else\textsubscript{\ensuremath{\beta}}\fi}
\newunicodechar{ₓ}{\ifmmode_{x}\else\textsubscript{\ensuremath{x}}\fi}
\newfontfamily\popdisplay{ArchivoBlack-Regular.ttf}[Path=../../../fonts/]
\newfontfamily\poplogo{SpaceGrotesk-Bold.ttf}[Path=../../../fonts/]
\newfontfamily\popsemi{PlusJakartaSans-SemiBold.ttf}[Path=../../../fonts/]
\newfontfamily\popxbold{PlusJakartaSans-ExtraBold.ttf}[Path=../../../fonts/]
\newfontfamily\popxboldls{PlusJakartaSans-ExtraBold.ttf}[Path=../../../fonts/,LetterSpace=3.0]

\setlist[enumerate]{leftmargin=1.7em,itemsep=5pt,topsep=4pt}
\setlist[itemize]{leftmargin=1.7em,itemsep=4pt,topsep=4pt,label={\color{poporange}\textbullet}}
\setlength{\parindent}{0pt}
\setlength{\parskip}{5pt}
\renewcommand{\arraystretch}{1.25}

% pgfplots : style Pop par défaut
\pgfplotsset{
  every axis/.append style={axis line style={popink,line width=1pt},tick style={popink},
    grid style={popline,line width=0.5pt},label style={font=\small},tick label style={font=\footnotesize}},
  popcurve/.style={poporange,line width=1.8pt,no markers,smooth},
}
% Même style disponible dans un \draw TikZ simple (hors pgfplots)
\tikzset{popcurve/.style={poporange,line width=1.8pt}}
\tikzset{popgrid/.style={popline,very thin}}
\colorlet{popcurve}{poporange} % le nom du style est parfois utilisé comme couleur

% ============================================================
% Pill / squiggle
% ============================================================
\newcommand{\poppill}[3]{%
  \tikz[baseline=-0.55ex]\node[draw=popink,line width=1.2pt,rounded corners=2.6pt,%
    fill=#2,inner xsep=6.5pt,inner ysep=3pt]{%
    \popxboldls\fontsize{7}{7}\selectfont\color{#3}\MakeUppercase{#1}};%
}
\newcommand{\popsquiggle}[1]{%
  \tikz[baseline=-0.4ex]{
    \draw[#1,line width=2.6pt,line cap=round]
      plot [smooth] coordinates {(0,0.09) (0.34,0.01) (0.68,0.21) (1.02,0.01) (1.36,0.10)};
  }%
}
% Mot-clé surligné (marqueur orange sous le texte)
\newcommand{\cle}[1]{{\popxbold\color{popdark}#1}}

% ============================================================
% En-tête / pied
% ============================================================
\pagestyle{fancy}
\fancyhf{}
\renewcommand{\headrulewidth}{0pt}
\renewcommand{\footrulewidth}{0pt}
\lhead{\raisebox{-0.15ex}{\poplogo\fontsize{13}{13}\selectfont\color{popink}OnBuch\color{poporange}+}}
\rhead{\poppill{\DOCMATIERE\ \textbullet\ \DOCNIVEAU\ \textbullet\ Leçon \DOCLECON}{white}{popink}}
\lfoot{\popxbold\fontsize{7.6}{7.6}\selectfont\color{popmuted}\MakeUppercase{Cours signé OnBuch+}\ \textbullet\ {\popsemi\DOCTITRE}}
\rfoot{\tikz[baseline=-0.6ex]\node[rounded corners=8.5pt,fill=popink,minimum height=17pt,minimum width=17pt,inner xsep=5pt,inner sep=0]{\popxbold\fontsize{8}{8}\selectfont\color{popcream}\thepage\,/\,\pageref*{LastPage}};}

% ============================================================
% Titres
% ============================================================
\newcommand{\chaptitre}[2]{%
  \begin{tcolorbox}[enhanced,colback=poporange,colframe=popink,boxrule=2.6pt,arc=11pt,
      drop shadow=popink,left=15pt,right=15pt,top=12pt,bottom=11pt,before skip=3pt,after skip=18pt]
    \poppill{#2}{popink}{popcream}\par\vspace{9pt}
    {\raggedright\hyphenpenalty=10000\exhyphenpenalty=10000\popdisplay\fontsize{22}{24}\selectfont\color{popcream}\MakeUppercase{#1}\par}
  \end{tcolorbox}
}
% Section numérotée : pastille orange avec le numéro + titre Archivo
\newcounter{popsec}
\newcommand{\coursec}[1]{%
  \stepcounter{popsec}\setcounter{popsub}{0}%
  \par\vspace{16pt}\noindent
  \begin{minipage}{\linewidth}
  \tikz[baseline=-0.7ex]\node[draw=popink,line width=1.6pt,fill=poporange,rounded corners=4pt,minimum size=22pt,inner sep=0]{\popdisplay\fontsize{12}{12}\selectfont\color{popcream}\thepopsec};%
  \hspace{8pt}{\popdisplay\fontsize{15}{17}\selectfont\color{popink}\MakeUppercase{#1}}\par
  \vspace{4pt}\noindent{\color{poporange}\rule{36pt}{3.6pt}}
  \end{minipage}\par\nopagebreak\vspace{8pt}
}
\newcounter{popsub}
\newcommand{\courssub}[1]{%
  \stepcounter{popsub}%
  \par\vspace{9pt}\noindent{\popxbold\fontsize{11.5}{13}\selectfont\color{popink}{\color{poporange}\thepopsec.\thepopsub}\hspace{6pt}#1}\par\nopagebreak\vspace{4pt}
}

% ============================================================
% Boîtes pédagogiques — même recette : fond blanc, bordure épaisse colorée,
% ombre dure encre, badge plein. Titre optionnel [#1].
% ============================================================
\tcbset{popbase/.style={breakable,enhanced,colback=white,boxrule=2.3pt,arc=9pt,
  shadow={3pt}{-3pt}{0pt}{popink},left=14pt,right=14pt,top=11pt,bottom=11pt,before skip=11pt,after skip=15pt,
  fontupper=\popsemi\fontsize{10.6}{14.5}\selectfont\color{popink},
  fontlower=\popsemi\fontsize{10.6}{14.5}\selectfont\color{popink}}}
% Coche dessinée (le glyphe ✓ n'existe pas dans les polices du document)
\newcommand{\popcheck}[1]{\tikz[baseline=-0.5ex]\draw[#1,line width=1.6pt,line cap=round,line join=round](-0.13,0.0)--(-0.04,-0.09)--(0.13,0.1);}
\providecommand{\coche}{\popcheck{popgreen}}
\providecommand{\resultat}[1]{\par\smallskip\noindent\fcolorbox{popgreen}{popgreenL}{\parbox{\dimexpr\linewidth-2\fboxsep-2\fboxrule\relax}{\textbf{Résultat.} #1}}\par\smallskip}
\newcommand{\popboxhead}[3]{\poppill{#1}{#2}{white}\ifx\relax#3\relax\else\hspace{6pt}{\popxbold\fontsize{10.6}{12}\selectfont\color{popink}#3}\fi\par\vspace{7pt}}

\newtcolorbox{aretenir}[1][]{popbase,colframe=popgreen,before upper={\popboxhead{À retenir}{popgreen}{#1}}}
% Texte en anglais (document de lecture, dialogue, extrait) : cadre
% d'étude. Un titre optionnel (source, auteur) s'affiche dans l'en-tête.
\newtcolorbox{textebox}[1][]{popbase,colframe=popblue,before upper={\popboxhead{Text}{popblue}{#1}\begin{otherlanguage}{english}},after upper={\end{otherlanguage}}}
% Marques ✓ / ✗ (forme correcte / fautive) souvent utilisées par les modèles
\providecommand{\cross}{\textcolor{poppink}{\ensuremath{\times}}}
\providecommand{\xmark}{\cross}\providecommand{\crossmark}{\cross}\providecommand{\cmark}{\ensuremath{\checkmark}}
% Ligne de réponse / justification à compléter (inventée par les modèles)
\providecommand{\justification}[1]{\par\noindent\textit{Justification~:} \dotfill\par}
% \textipa (package tipa non chargé) : la police ne couvre pas l'API -> simple passage
\providecommand{\textipa}[1]{#1}
% Soulignement (ulem non chargé) : \uline, \ul
\providecommand{\uline}[1]{\underline{#1}}\providecommand{\ul}[1]{\underline{#1}}
% \glossaire{\item ...} inventée par les modèles : liste de glossaire
\providecommand{\glossaire}[1]{\begin{itemize}#1\end{itemize}}
% \alert (beamer) inventée par les modèles : texte mis en évidence
\providecommand{\alert}[1]{\textbf{\textcolor{poppink}{#1}}}
% variantes ASCII d'ordinaux écrites en macro
\providecommand{\iere}{\textsuperscript{re}}\providecommand{\ieme}{\textsuperscript{e}}\providecommand{\ere}{\textsuperscript{re}}\providecommand{\eme}{\textsuperscript{e}}\providecommand{\er}{\textsuperscript{er}}\providecommand{\ie}{\textsuperscript{e}}
% Ordinaux français écrits en macro par les modèles : 1\ière, 3\ième
\newcommand{\ière}{\textsuperscript{re}}\newcommand{\ième}{\textsuperscript{e}}
% \exemple (inventée par les modèles) : étiquette « Exemple : »
\providecommand{\exemple}{\textbf{Exemple~:}\ }
% \square utilisable aussi hors mode math (cases à cocher)
\AtBeginDocument{\let\squareMath\square\renewcommand{\square}{\ensuremath{\squareMath}}}
% Mot ou phrase en anglais à l'intérieur d'un paragraphe français
\newcommand{\eng}[1]{\ifmmode\text{\textit{#1}}\else\begingroup\selectlanguage{english}\itshape #1\endgroup\fi}
\let\esp\eng % alias toléré
\newtcolorbox{exemplebox}[1][]{popbase,colframe=poporange,before upper={\popboxhead{Exemple}{poporange}{#1}}}
\newtcolorbox{definition}[1][]{popbase,colframe=popblue,before upper={\popboxhead{Définition}{popblue}{#1}}}
\newtcolorbox{propriete}[1][]{popbase,colframe=poppurple,before upper={\popboxhead{Propriété}{poppurple}{#1}}}
\newtcolorbox{methode}[1][]{popbase,colframe=popgold,before upper={\popboxhead{Méthode}{popgold}{#1}}}
\newtcolorbox{attention}[1][]{popbase,colframe=poppink,before upper={\popboxhead{Attention}{poppink}{#1}}}
\newtcolorbox{experience}[1][]{popbase,colframe=popblue,colback=popblueL,before upper={\popboxhead{Expérience}{popblue}{#1}}}
% « Le savais-tu ? » : carte violette pleine (contexte, culture, Cameroun)
\newtcolorbox{savaistu}[1][]{popbase,colback=poppurple,colframe=popink,
  fontupper=\popsemi\fontsize{10.4}{14.2}\selectfont\color{white},
  before upper={\poppill{Le savais-tu\,?}{white}{poppurple}\ifx\relax#1\relax\else\hspace{6pt}{\popxbold\color{white}#1}\fi\par\vspace{7pt}}}
% Exercice résolu : énoncé en haut, \tcblower puis la solution sur fond crème
\newcounter{popexo}\newcounter{popexr}
\newtcolorbox{exoresolu}[1][]{popbase,colframe=poporange,colbacklower=popcream,
  segmentation style={popink,line width=1.2pt,dashed},
  before upper={\stepcounter{popexr}\popboxhead{Exercice résolu \thepopexr}{poporange}{#1}},
  before lower={\poppill{Solution}{popink}{popcream}\par\vspace{6pt}}}
% Exercice (non résolu en place) — la correction est dans la partie Corrigés
\newtcolorbox{exercice}[1][]{popbase,colframe=popink,
  before upper={\stepcounter{popexo}\popboxhead{Exercice \thepopexo}{popink}{#1}}}
% Corrigé
\newtcolorbox{corrige}[1][]{popbase,colframe=popgreen,colback=popgreenL,
  before upper={\popboxhead{Corrigé}{popgreen}{#1}}}
% Objectifs / prérequis / bilan
\newtcolorbox{objectifs}{popbase,colframe=popink,colback=poporangeL,
  before upper={\popboxhead{Objectifs de la leçon}{popink}{}}}
\newtcolorbox{prerequis}{popbase,colframe=popink,colback=popblueL,
  before upper={\popboxhead{Prérequis}{popblue}{}}}
\newtcolorbox{bilan}{popbase,colframe=popink,colback=white,boxrule=2.8pt,
  before upper={\popboxhead{Fiche bilan}{poporange}{L'essentiel à connaître pour l'examen}}}
% Figure encadrée avec légende
\newtcolorbox{popfigure}[1][]{enhanced,breakable=false,colback=white,colframe=popline,boxrule=1.4pt,arc=8pt,
  left=10pt,right=10pt,top=10pt,bottom=8pt,before skip=10pt,after skip=12pt,halign=center,
  lower separated=false,halign lower=center,
  fontlower=\popsemi\fontsize{9}{11.5}\selectfont\color{popmuted},
  #1}
\newcommand{\legende}[1]{\par\vspace{6pt}{\popsemi\fontsize{9}{11.5}\selectfont\color{popmuted}#1}}

% ============================================================
% Signature OnBuch+
% ============================================================
\newcommand{\poplogoinline}[1]{{\poplogo\fontsize{#1}{#1}\selectfont\color{popink}OnBuch\color{poporange}+}}
\newcommand{\signatureonbuch}{%
  \begin{tcolorbox}[enhanced,colback=popink,colframe=popink,boxrule=2.2pt,arc=10pt,
      drop shadow=poporange,left=14pt,right=14pt,top=10pt,bottom=10pt,before skip=0pt,after skip=0pt]
    \tikz[baseline=-0.7ex]\node[circle,fill=popgreen,draw=popcream,line width=1.3pt,minimum size=22pt,inner sep=0]{\popcheck{white}};%
    \hspace{9pt}\begin{minipage}[c]{0.62\linewidth}
      {\popxbold\fontsize{12}{13}\selectfont\color{popcream}Signé\ \textbullet\ {\poplogo OnBuch\color{poporange}+}}\par\vspace{2pt}
      {\raggedright\popsemi\fontsize{8}{10.5}\selectfont\color{popline}\MakeUppercase{Équipe pédagogique OnBuch+}\\\MakeUppercase{Conforme au programme officiel MINESEC}\par}
    \end{minipage}\hfill
    {\poplogo\fontsize{22}{22}\selectfont\color{popcream}OnBuch\color{poporange}+}
  \end{tcolorbox}%
}

% ============================================================
% Page de garde
% #1 titre · #2 matière · #3 niveau · #4 module · #5 n° leçon · #6 durée ·
% #7 description / objectifs (texte ou itemize)
% ============================================================
\newcommand{\pagegarde}[7]{%
  \thispagestyle{empty}
  \newgeometry{a4paper,margin=1.7cm,top=1.6cm,bottom=1.4cm}%
  \begin{tikzpicture}[remember picture,overlay]
    \fill[poporange!18] ([xshift=-3.2cm,yshift=-3.6cm]current page.north east) circle (4.6cm);
    \fill[poppurple!14] ([xshift=2.4cm,yshift=7.6cm]current page.south west) circle (3.8cm);
  \end{tikzpicture}%
  {\poplogo\fontsize{30}{30}\selectfont\color{popink}OnBuch\color{poporange}+}\hfill
  \raisebox{0.6ex}{\poppill{Cours signé \textbullet\ Premium}{popink}{popcream}}\par
  \vspace{1pt}\popsquiggle{poppurple}\par
  \vspace{8pt}
  \begin{tcolorbox}[enhanced,colback=poporange,colframe=popink,boxrule=2.8pt,arc=13pt,
      drop shadow=popink,left=18pt,right=18pt,top=12pt,bottom=13pt,after skip=10pt]
    \poppill{#2 \textbullet\ #3}{popink}{popcream}\hspace{5pt}\poppill{#4}{white}{popink}\par\vspace{8pt}
    {\popdisplay\fontsize{14}{14}\selectfont\color{popink}LEÇON #5}\par\vspace{4pt}
    {\raggedright\hyphenpenalty=10000\exhyphenpenalty=10000\popdisplay\fontsize{25}{27}\selectfont\color{popcream}\MakeUppercase{#1}\par}
    \vspace{8pt}
    {\popxbold\fontsize{9.5}{9.5}\selectfont\color{popink}\MakeUppercase{Durée conseillée : #6}}
  \end{tcolorbox}
  \begin{tcolorbox}[enhanced,colback=white,colframe=popink,boxrule=2.2pt,arc=10pt,
      drop shadow=popink,left=16pt,right=16pt,top=10pt,bottom=10pt,after skip=8pt]
    \poppill{Dans cette leçon}{poppurple}{white}\par\vspace{5pt}
    \popsemi\fontsize{9.3}{12.6}\selectfont\color{popink}#7
  \end{tcolorbox}
  \noindent\begin{minipage}[t]{0.487\linewidth}
    \begin{tcolorbox}[enhanced,colback=popgreen,colframe=popink,boxrule=2.2pt,arc=10pt,
        drop shadow=popink,left=11pt,right=11pt,top=10pt,bottom=10pt,height=3.1cm,valign=center]
      \tcbox[colback=white,colframe=popink,boxrule=1.3pt,arc=5pt,boxsep=0pt,nobeforeafter,
        left=3pt,right=3pt,top=3pt,bottom=3pt,valign=center]{\qrcode[height=1.9cm]{https://play.google.com/store/apps/details?id=com.luvvix.onbuch&hl=fr}}%
      \hspace{8pt}\begin{minipage}[c]{0.5\linewidth}
        {\popdisplay\fontsize{11}{12.5}\selectfont\color{white}\MakeUppercase{Télécharge OnBuch}}\par\vspace{4pt}
        {\raggedright\popsemi\fontsize{8.4}{11}\selectfont\color{white}Gratuit \textbullet\ Google Play\\Tuteur IA, annales, concours, résultats\par}
      \end{minipage}
    \end{tcolorbox}
  \end{minipage}\hfill
  \begin{minipage}[t]{0.487\linewidth}
    \begin{tcolorbox}[enhanced,colback=poppurple,colframe=popink,boxrule=2.2pt,arc=10pt,
        drop shadow=popink,left=11pt,right=11pt,top=10pt,bottom=10pt,height=3.1cm,valign=center]
      \tikz[baseline=-0.7ex]\node[circle,fill=white,draw=popink,line width=1.3pt,minimum size=1.6cm,inner sep=0]{\popdisplay\fontsize{24}{24}\selectfont\color{poppurple}+};%
      \hspace{8pt}\begin{minipage}[c]{0.6\linewidth}
        {\popdisplay\fontsize{11}{12.5}\selectfont\color{white}\MakeUppercase{Passe à OnBuch+}}\par\vspace{4pt}
        {\raggedright\popsemi\fontsize{8.4}{11}\selectfont\color{white}Tous les cours signés, Tuteur IA illimité, fiches PDF — onglet OnBuch+ de l'app\par}
      \end{minipage}
    \end{tcolorbox}
  \end{minipage}
  \vspace{8pt}
  \popcontenu
  \vfill
  \signatureonbuch
  \restoregeometry
  \newpage
}

% Rangée « Ce cours contient » (page de garde)
\newcommand{\poptile}[3]{% #1 couleur #2 gros texte #3 légende
  \begin{tcolorbox}[enhanced,colback=#1,colframe=popink,boxrule=2pt,arc=9pt,shadow={2.5pt}{-2.5pt}{0pt}{popink},
      left=6pt,right=6pt,top=9pt,bottom=9pt,halign=center,valign=center,height=1.75cm,width=\linewidth,nobeforeafter]
    {\popdisplay\fontsize{12.5}{13}\selectfont\color{white}\MakeUppercase{#2}}\par\vspace{3pt}
    {\centering\popsemi\fontsize{7.8}{9.5}\selectfont\color{white}#3\par}
  \end{tcolorbox}}
\newcommand{\popcontenu}{%
  \par\noindent{\popxboldls\fontsize{7.5}{7.5}\selectfont\color{popmuted}CE COURS SIGNÉ CONTIENT}\par\vspace{7pt}
  \noindent\begin{minipage}[t]{0.235\linewidth}\poptile{popblue}{Cours}{complet, illustré, pas à pas}\end{minipage}\hfill
  \begin{minipage}[t]{0.235\linewidth}\poptile{poporange}{Méthodes}{et exercices résolus}\end{minipage}\hfill
  \begin{minipage}[t]{0.235\linewidth}\poptile{poppink}{Exercices}{type examen + corrigés}\end{minipage}\hfill
  \begin{minipage}[t]{0.235\linewidth}\poptile{popgreen}{Fiche bilan}{l'essentiel à retenir}\end{minipage}\par}

% Sommaire
\newtcolorbox{sommaire}{enhanced,colback=white,colframe=popink,boxrule=2.2pt,arc=10pt,
  drop shadow=popink,left=18pt,right=18pt,top=13pt,bottom=13pt,before skip=6pt,after skip=20pt,
  before upper={\poppill{Sommaire}{poppurple}{white}\par\vspace{9pt}\popsemi\fontsize{10.6}{14.5}\selectfont\color{popink}}}

% Signature de fin de cours — poussée tout en bas de la dernière page
\newcommand{\findecours}{%
  \par\vfill\nopagebreak
  \begin{center}\popsquiggle{poporange}\end{center}\vspace{4pt}
  \signatureonbuch
}
\AtBeginDocument{\def\llbracket{\mathopen{[\mkern-3.5mu[}}\def\rrbracket{\mathclose{]\mkern-3.5mu]}}} % intervalles d'entiers (glyphe ⟦ absent de la police)
% Glyphes absents de Fira Math : redéfinis à partir de glyphes disponibles
\AtBeginDocument{%
  \def\setminus{\mathbin{\backslash}}\let\smallsetminus\setminus
  \def\vdots{\mathord{\vbox{\baselineskip3.2pt\lineskiplimit0pt\kern4pt\hbox{.}\hbox{.}\hbox{.}}}}%
  \def\ddots{\mathinner{\mkern1mu\raise6.5pt\hbox{.}\mkern2mu\raise3.6pt\hbox{.}\mkern2mu\raise0.7pt\hbox{.}\mkern1mu}}%
  \def\triangle{\mathord{\scalebox{0.9}{$\symup{\Delta}$}}}%
  \def\nabla{\mathord{\raisebox{\depth}{\scalebox{1}[-1]{$\symup{\Delta}$}}}}%
  \def\checkmark{\popcheck{popdark}}%
  \def\star{\mathbin{\ast}}\def\bigstar{\mathord{\ast}}%
  \def\diamond{\mathbin{\scalebox{0.6}{$\Diamond$}}}%
  \def\bigcup{\mathop{\mathchoice{\raisebox{-0.3em}{\scalebox{1.7}{$\cup$}}}{\scalebox{1.25}{$\cup$}}{\cup}{\cup}}\displaylimits}%
  \def\bigcap{\mathop{\mathchoice{\raisebox{-0.3em}{\scalebox{1.7}{$\cap$}}}{\scalebox{1.25}{$\cap$}}{\cap}{\cap}}\displaylimits}%
  \def\lesseqqgtr{\mathrel{\lessgtr}}\def\bigoplus{\mathop{\oplus}\displaylimits}%
}
\providecommand{\ssi}{si et seulement si} % abréviation fréquente
\AtBeginDocument{\providecommand{\wideparen}{\overparen}} % arc de cercle

%% FIGFIX-BEGIN v5 — correctifs globaux des figures (docs/onbuch-generation/tools/figfix)
\usepackage{adjustbox}\usepackage{etoolbox}\tcbuselibrary{raster}
% toute figure TikZ/pgfplots plus large que la ligne est réduite à la largeur disponible
\BeforeBeginEnvironment{tikzpicture}{\begin{adjustbox}{max width=\linewidth}}
\AfterEndEnvironment{tikzpicture}{\end{adjustbox}}
% légendes pgfplots lisibles par-dessus les courbes
\pgfplotsset{every axis/.append style={legend style={fill=white,fill opacity=0.92,text opacity=1,draw=black!15,inner sep=3pt}}}
% étiquettes d'arêtes (syntaxe edge["..."]) avec fond blanc
\tikzset{every edge quotes/.append style={fill=white,inner sep=1.2pt}}
%% FIGFIX-END
\begin{document}

\pagegarde{Lesson 7 — Reading: Texts on recreational activities}{Anglais}{Tle A}{Chapter 1 : Family and Social Life}{EN07}{2 h}{Cette leçon de compréhension écrite propose des textes authentiques sur les loisirs et activités récréatives dans le monde anglophone. Elle entraîne les élèves aux stratégies de lecture (skimming, scanning, déduction du sens) et au lexique du thème, en vue de l'épreuve de reading comprehension du Baccalauréat A.
\vspace{4pt}
\begin{multicols}{2}\small
\begin{itemize}
\item À la fin de cette leçon, je sais identifier l'idée générale d'un texte sur les loisirs grâce au skimming
\item repérer des informations précises (dates, lieux, personnes, chiffres) grâce au scanning
\item déduire le sens de mots inconnus à partir du contexte
\item utiliser le lexique des activités récréatives (sports, jeux, arts, médias) en anglais
\item répondre à des questions de compréhension de types variés (vrai/faux, QCM, réponses courtes)
\item distinguer les faits des opinions dans un texte
\end{itemize}\end{multicols}}

\begin{sommaire}
\begin{multicols}{2}
\begin{enumerate}
\item Warm-up: Talking about free time
\item Reading text: 'How We Spend Our Free Time'
\item Reading strategies: skimming, scanning, guessing
\item Comprehension practice
\item Vocabulary: the world of recreation
\item Guided production: summary and opinion
\item Activité d'intégration
\item Méthodes et exercices résolus
\item Exercices
\item Corrigés des exercices
\item Fiche bilan
\end{enumerate}
\end{multicols}
\end{sommaire}

\begin{objectifs}
\begin{itemize}
\item À la fin de cette leçon, je sais identifier l'idée générale d'un texte sur les loisirs grâce au skimming
\item repérer des informations précises (dates, lieux, personnes, chiffres) grâce au scanning
\item déduire le sens de mots inconnus à partir du contexte
\item utiliser le lexique des activités récréatives (sports, jeux, arts, médias) en anglais
\item répondre à des questions de compréhension de types variés (vrai/faux, QCM, réponses courtes)
\item distinguer les faits des opinions dans un texte
\item résumer un texte en quelques phrases avec mes propres mots
\item exprimer mon opinion sur une activité de loisir en anglais
\end{itemize}
\end{objectifs}

\begin{prerequis}
\begin{itemize}
\item Le présent simple et le present continuous (Première)
\item Le vocabulaire de base des sports et des hobbies (Seconde/Première)
\item Les verbes de goût : like, enjoy, prefer, hate + -ing
\item Les connecteurs simples : and, but, because, however
\item Les wh-questions (who, what, where, when, why, how)
\end{itemize}
\end{prerequis}

\newpage

\coursec{Warm-up: Talking about free time}

\courssub{Situation d'accroche : un samedi après-midi à travers le monde anglophone}

\begin{textebox}[Saturday afternoon around the English-speaking world]
It is Saturday afternoon in Bamenda. Some students are playing football on the field, others are watching a Nollywood film at home, and a group of friends is rehearsing traditional dances for the school cultural day. In London, teenagers of the same age are playing video games online, while in Lagos others are attending a music concert. In Buea, a church choir is practising gospel songs for Sunday service, and in a village near Bamenda, a farmer is resting under a mango tree, listening to the radio. Everyone is doing something different — but nobody is working. All these people are enjoying their \emph{free time}.
\end{textebox}

\textbf{Problématique.} How do people around the English-speaking world spend their free time? What exactly do we mean by \eng{recreation}, \eng{leisure} or \eng{hobby}? Today, we will read texts about \cle{recreational activities} and learn how to extract information quickly and efficiently — an essential skill for the Bac. But before opening any text, a good reader always \emph{warms up}: he or she activates what is already known about the topic. That is exactly what we are going to do now.

\begin{savaistu}[Leisure in Anglophone Cameroon]
In the English-speaking regions of Cameroon, \eng{choral singing} (le chant choral) and \eng{traditional dance groups} are extremely popular \emph{community} leisure activities: almost every church, school and village has its choir or dance group, and rehearsals are important social events. In the United Kingdom, the equivalent community passion is \emph{football}: millions of people play in local clubs or watch matches every weekend. In both cultures, recreation is not only about fun — it builds friendship and community spirit.
\end{savaistu}

\courssub{Brainstorming : what do YOU do in your free time?}

Avant de lire un texte, interrogeons notre propre expérience. Regardez la situation d'accroche et pensez à votre propre samedi après-midi. Quelles activités pratiquez-vous quand vous ne travaillez pas ?

\begin{experience}[Class brainstorming : “Our free time”]
\begin{enumerate}
\item En deux minutes, chaque élève écrit en anglais \textbf{trois} activités qu'il ou elle pratique pendant son temps libre. Modèle de phrase : \eng{In my free time, I play football.} / \eng{I enjoy dancing.} / \eng{I like reading novels.}
\item Le professeur recueille les réponses au tableau et les classe en colonnes : \eng{sports}, \eng{arts and music}, \eng{games}, \eng{media}, \eng{outdoor activities}.
\item Questions de relance en anglais : \eng{How often do you do it?} (every day, twice a week, at weekends) — \eng{Where do you do it?} (at home, on the field, at church) — \eng{Who do you do it with?} (alone, with friends, with my family).
\item Réponses attendues typiques au Cameroun : \eng{playing football}, \eng{dancing}, \eng{singing in the church choir}, \eng{reading}, \eng{playing video games}, \eng{watching films}, \eng{helping on the family farm}, \eng{playing draughts or ludo}, \eng{chatting with friends}, \eng{listening to music}.
\end{enumerate}
\end{experience}

\begin{propriete}[Verbes de préférence suivis du gérondif (-ing)]
Les verbes \eng{like}, \eng{love}, \eng{enjoy}, \eng{hate}, \eng{prefer}, \eng{mind}, \eng{avoid}, \eng{finish}, \eng{suggest} sont suivis du \textbf{gérondif} (forme en \textbf{-ing}), jamais de l'infinitif avec \emph{to} (sauf \eng{like} / \eng{prefer} / \eng{love} / \eng{hate} qui acceptent les deux formes avec une nuance de sens).
\begin{center}
\begin{tabular}{|l|l|l|}
\hline
\rowcolor{popblueL} \textbf{Structure} & \textbf{Exemple correct} & \textbf{Traduction} \\
\hline
\eng{enjoy + V-ing} & \eng{I enjoy \textbf{dancing}.} & « J'aime danser / Je prends du plaisir à danser. » \\
\hline
\eng{like + V-ing} (habitude, goût général) & \eng{I like \textbf{reading} novels.} & « J'aime lire des romans. » \\
\hline
\eng{like + to V} (choix ponctuel, utilité) & \eng{I like \textbf{to read} before sleeping.} & « J'aime lire avant de dormir. » \\
\hline
\end{tabular}
\end{center}
\end{propriete}

\begin{attention}[Attention aux structures après “like” et “enjoy” — Pièges fréquents]
\begin{itemize}
\item Après \eng{enjoy}, \textbf{jamais} d'infinitif : \eng{*I enjoy to dance} $\rightarrow$ \textbf{FAUX}. Dire \eng{I enjoy \textbf{dancing}}.
\item Après \eng{like} / \eng{love} / \eng{hate} / \eng{prefer}, les deux formes existent mais \textbf{préférez le -ing} pour exprimer un loisir, une habitude : \eng{I like swimming} (plus naturel que \eng{I like to swim} pour un hobby).
\item \textbf{Play + sport} : sans article. \eng{play football}, \eng{play basketball}, \eng{play handball}.
\item \textbf{Play + instrument de musique} : avec \textbf{the}. \eng{play \textbf{the} piano}, \eng{play \textbf{the} guitar}, \eng{play \textbf{the} drums}.
\end{itemize}
\end{attention}

\courssub{Leisure, recreation, hobby, pastime : quatre mots à ne pas confondre}

Le français utilise surtout le mot « loisir(s) » ; l'anglais possède quatre mots proches mais distincts. Observons d'abord ces exemples tirés de situations réelles :

\begin{itemize}
\item \eng{What do you do in your \textbf{leisure time}?} — « Que fais-tu pendant ton temps libre ? »
\item \eng{The town council has built a new \textbf{recreation} centre.} — « La municipalité a construit un nouveau centre de loisirs. »
\item \eng{Her favourite \textbf{hobby} is collecting stamps.} — « Son passe-temps favori est la philatélie. »
\item \eng{Reading is a relaxing \textbf{pastime}.} — « La lecture est un passe-temps relaxant. »
\end{itemize}

\begin{definition}[Recreation]
\cle{Recreation} (prononciation : « ré-kri-é-cheune ») : \emph{activities people do for enjoyment when they are not working} — les activités que les gens pratiquent pour le plaisir quand ils ne travaillent pas. En français : « loisirs », « activités récréatives ». L'adjectif correspondant est \eng{recreational} : \eng{recreational activities} = « activités de loisir ».
\end{definition}

\begin{center}
\begin{tabularx}{\linewidth}{|l|X|X|X|}
\hline
\rowcolor{popblueL} \textbf{Mot anglais} & \textbf{Sens précis} & \textbf{Traduction française} & \textbf{Exemple} \\
\hline
\eng{free time / leisure time} & le temps dont on dispose hors travail et études & temps libre & \eng{I read in my free time.} \\
\hline
\eng{leisure} & notion générale et abstraite (souvent invariable) & loisir(s) (notion) & \eng{Leisure is important for health.} \\
\hline
\eng{recreation} & les activités pratiquées pour se détendre et s'amuser & loisirs, récréation & \eng{Swimming is a popular recreation.} \\
\hline
\eng{hobby} & une activité régulière, souvent solitaire et pratiquée avec passion (collection, couture, jardinage...) & passe-temps favori, hobby & \eng{Gardening is his hobby.} \\
\hline
\eng{pastime} & une activité qui fait passer le temps agréablement & passe-temps & \eng{Draughts is a common pastime in Cameroon.} \\
\hline
\end{tabularx}
\end{center}

\begin{attention}[Faux amis et pièges de traduction]
\begin{itemize}
\item \eng{Recreation} ne signifie pas « récréation » au sens de la pause entre les cours : la récréation scolaire se dit \eng{break} ou \eng{break time}.
\item Ne traduisez pas « mes loisirs » par \emph{my leisures} : \eng{leisure} est en général \textbf{invariable}. Dites \eng{my leisure activities} ou \eng{my hobbies}.
\item Pluriel : \eng{hobby} devient \eng{hobbies} (y $\rightarrow$ ies), comme \eng{activity} $\rightarrow$ \eng{activities}.
\end{itemize}
\end{attention}

\courssub{Classifying activities : indoor vs outdoor, active vs passive}

Pour parler des loisirs avec précision, les Anglophones classent les activités selon deux grands critères.

\begin{definition}[Deux classifications essentielles]
\begin{itemize}
\item \cle{Indoor activities} (activités d'intérieur) : pratiquées à l'intérieur — \eng{reading, playing video games, watching television, playing board games, cooking}.
\item \cle{Outdoor activities} (activités de plein air) : pratiquées dehors — \eng{playing football, farming, hiking, cycling, fishing}.
\item \cle{Active recreation} (loisirs actifs) : le corps ou l'esprit produit un effort — \eng{dancing, singing in a choir, playing sports, gardening, reading, playing video games}.
\item \cle{Passive recreation} (loisirs passifs) : on reçoit sans produire d'effort conscient — \eng{watching films, listening to music, scrolling on the phone}.
\end{itemize}
\end{definition}

\begin{center}
\begin{tabularx}{\linewidth}{|X|X|X|}
\hline
\rowcolor{popblueL} \textbf{Activity} & \textbf{Indoor / Outdoor} & \textbf{Active / Passive} \\
\hline
\eng{playing football} & outdoor & active \\
\hline
\eng{watching a Nollywood film} & indoor & passive \\
\hline
\eng{rehearsing traditional dances} & indoor or outdoor & active \\
\hline
\eng{reading a novel} & indoor & active (mental effort) \\
\hline
\eng{listening to the radio} & indoor or outdoor & passive \\
\hline
\eng{playing video games} & indoor & active (mental effort) \\
\hline
\eng{working on the family farm} & outdoor & active \\
\hline
\end{tabularx}
\end{center}

\begin{popfigure}
\begin{center}\tcbox[colback=popblue,colframe=popblue,coltext=white,arc=9pt,boxsep=4pt,left=6pt,right=6pt]{\bfseries\small Recreational activities}\end{center}
\vspace{-2pt}
\begin{tcbraster}[raster columns=3,raster equal height=rows,raster column skip=6pt,raster row skip=6pt,enhanced,blankest]
\begin{tcolorbox}[enhanced,colback=white,colframe=poporange,boxrule=0.9pt,arc=3pt,title={Sports},fonttitle=\bfseries\scriptsize,coltitle=white,colbacktitle=poporange,left=4pt,right=4pt,top=2pt,bottom=2pt,boxsep=1pt,toptitle=1.5pt,bottomtitle=1.5pt,halign=left]
\begin{itemize}[nosep,leftmargin=*,itemsep=1pt,font=\scriptsize]
\item football
\item athletics
\end{itemize}
\end{tcolorbox}
\begin{tcolorbox}[enhanced,colback=white,colframe=poppurple,boxrule=0.9pt,arc=3pt,title={Arts \& Music},fonttitle=\bfseries\scriptsize,coltitle=white,colbacktitle=poppurple,left=4pt,right=4pt,top=2pt,bottom=2pt,boxsep=1pt,toptitle=1.5pt,bottomtitle=1.5pt,halign=left]
\begin{itemize}[nosep,leftmargin=*,itemsep=1pt,font=\scriptsize]
\item choir singing
\item dancing
\end{itemize}
\end{tcolorbox}
\begin{tcolorbox}[enhanced,colback=white,colframe=popgreen,boxrule=0.9pt,arc=3pt,title={Games},fonttitle=\bfseries\scriptsize,coltitle=white,colbacktitle=popgreen,left=4pt,right=4pt,top=2pt,bottom=2pt,boxsep=1pt,toptitle=1.5pt,bottomtitle=1.5pt,halign=left]
\begin{itemize}[nosep,leftmargin=*,itemsep=1pt,font=\scriptsize]
\item video games
\item draughts, ludo
\end{itemize}
\end{tcolorbox}
\begin{tcolorbox}[enhanced,colback=white,colframe=poppink,boxrule=0.9pt,arc=3pt,title={Media},fonttitle=\bfseries\scriptsize,coltitle=white,colbacktitle=poppink,left=4pt,right=4pt,top=2pt,bottom=2pt,boxsep=1pt,toptitle=1.5pt,bottomtitle=1.5pt,halign=left]
\begin{itemize}[nosep,leftmargin=*,itemsep=1pt,font=\scriptsize]
\item films
\item radio, music
\end{itemize}
\end{tcolorbox}
\begin{tcolorbox}[enhanced,colback=white,colframe=popgold!80!black,boxrule=0.9pt,arc=3pt,title={Outdoor activities},fonttitle=\bfseries\scriptsize,coltitle=white,colbacktitle=popgold!80!black,left=4pt,right=4pt,top=2pt,bottom=2pt,boxsep=1pt,toptitle=1.5pt,bottomtitle=1.5pt,halign=left]
\begin{itemize}[nosep,leftmargin=*,itemsep=1pt,font=\scriptsize]
\item hiking
\item farming, fishing
\end{itemize}
\end{tcolorbox}
\end{tcbraster}

\legende{Carte mentale du lexique : le monde des activités récréatives.}
\end{popfigure}

\begin{exoresolu}[Classify the activities]
\textbf{Énoncé.} Put each activity in the correct column: \emph{indoor} or \emph{outdoor}. Activities: \eng{playing video games — hiking — watching television — fishing — singing in a choir — playing football}.
\tcblower
\textbf{Solution détaillée.}
\begin{itemize}
\item \eng{playing video games} $\rightarrow$ \textbf{indoor} (on joue à la maison ou dans une salle de jeux).
\item \eng{hiking} $\rightarrow$ \textbf{outdoor} (la randonnée se pratique en pleine nature).
\item \eng{watching television} $\rightarrow$ \textbf{indoor}.
\item \eng{fishing} $\rightarrow$ \textbf{outdoor} (au bord d'une rivière ou d'un lac).
\item \eng{singing in a choir} $\rightarrow$ \textbf{indoor} (à l'église ou dans une salle) — mais une chorale peut aussi chanter dehors lors d'une fête ; le contexte décide.
\item \eng{playing football} $\rightarrow$ \textbf{outdoor} (sur un terrain).
\end{itemize}
\textbf{Méthode :} posez-vous la question \eng{Where do people usually do this activity?} — le mot \eng{usually} est important, car beaucoup d'activités peuvent se pratiquer dans les deux cadres.
\end{exoresolu}

\courssub{Predicting from the title : what will the text be about?}

Un bon lecteur ne commence jamais un texte « à froid ». Avant de lire la première ligne, il \emph{prédit} le contenu. Cette stratégie, appelée \cle{predicting}, est systématiquement évaluée au Bac dans les épreuves de compréhension de l'écrit.

\begin{methode}[How to predict the content of a text]
\begin{enumerate}
\item \textbf{Read the title carefully.} Soulignez les mots-clés du titre. Ici : \eng{How We Spend Our Free Time}. Mots-clés : \eng{how} (une manière de faire), \eng{we} (un groupe de personnes), \eng{spend} (employer, consacrer), \eng{free time} (temps libre).
\item \textbf{Look at the illustrations.} Les images qui accompagnent un texte donnent des indices sur les lieux, les personnages et les activités décrits.
\item \textbf{Ask yourself questions before reading.} Transformez le titre en questions : \eng{Who is “we”? Which activities will be mentioned? Will the text compare different countries or different age groups?}
\item \textbf{Activate your schema.} Rappelez-vous ce que vous savez déjà du sujet (votre propre expérience, le brainstorming de la classe). Votre cerveau lira plus vite car il sait à quoi s'attendre.
\item \textbf{Read the first sentence of each paragraph} (si le temps manque) pour confirmer ou corriger vos prédictions.
\end{enumerate}
\end{methode}

\begin{exemplebox}[Applying the method to our text]
Titre : \eng{How We Spend Our Free Time}. Prédictions possibles d'un élève de Tle A :
\begin{itemize}
\item \eng{The text will probably describe several leisure activities.} (Le texte décrira probablement plusieurs activités de loisir.)
\item \eng{It may compare young people and adults, or town and village life.} (Il comparera peut-être jeunes et adultes, ou ville et village.)
\item \eng{It may mention sports, music and media, because these are the most common free-time activities.} (Il mentionnera peut-être le sport, la musique et les médias.)
\end{itemize}
Ces hypothèses seront vérifiées pendant la lecture : c'est le principe du cycle \emph{predict — read — check}.
\end{exemplebox}

\courssub{Skimming and scanning : two essential reading techniques}

Une fois les prédictions posées, on entre dans le texte avec deux techniques complémentaires, explicitement demandées au Baccalauréat.

\begin{definition}[Skimming — Lecture rapide pour l'idée générale]
\cle{Skimming} consiste à lire \textbf{vite} (environ 200--300 mots/minute) en ne fixant que les mots porteurs de sens (noms, verbes, adjectifs) pour répondre à la question : \eng{What is the text about?} (De quoi parle le texte ?). On saute les mots inconnus, les détails, les exemples.
\end{definition}

\begin{definition}[Scanning — Lecture ciblée pour une information précise]
\cle{Scanning} consiste à \textbf{balayer} le texte des yeux à la recherche d'un mot-clé précis (un chiffre, un nom propre, une date, un mot technique) pour répondre à une question du type : \eng{Where does the text mention...?} / \eng{How many...?} / \eng{When...?} On ne lit pas le texte, on le scanne comme un code-barres.
\end{definition}

\begin{methode}[Skimming \& Scanning step by step]
\begin{enumerate}
\item \textbf{Skimming first :} Lisez le premier et le dernier paragraphe, puis la première phrase de chaque paragraphe intermédiaire. Notez l'idée principale en une phrase.
\item \textbf{Read the questions :} Identifiez le type d'information cherchée (chiffre $\rightarrow$ scanning ; idée générale $\rightarrow$ skimming ; opinion $\rightarrow$ lecture détaillée).
\item \textbf{Scanning for details :} Cherchez les mots-clés de la question dans le texte (noms propres, chiffres, mots en gras/italiques). Lisez \emph{autour} du mot trouvé.
\item \textbf{Guessing meaning :} Pour les mots inconnus nécessaires à la réponse, utilisez les indices de contexte (voir méthode ci-dessous).
\end{enumerate}
\end{methode}

\courssub{Guessing meaning from context : déduire le sens des mots inconnus}

Au Bac, vous ne pouvez pas utiliser de dictionnaire. La stratégie \cle{guessing meaning from context} (inférence lexicale) est donc vitale.

\begin{methode}[How to guess the meaning of an unknown word]
\begin{enumerate}
\item \textbf{Identify the part of speech.} Est-ce un nom, un verbe, un adjectif, un adverbe ? (Position dans la phrase, suffixes : \emph{-tion, -ment, -ness} = nom ; \emph{-ful, -ous, -ive} = adjectif ; \emph{-ly} = adverbe).
\item \textbf{Look for context clues (indices contextuels) :}
\begin{itemize}
\item \textbf{Definition / Explanation :} \eng{... recreation, \textbf{that is,} activities done for fun...} / \eng{... recreation, \textbf{i.e.,} hobbies...}
\item \textbf{Synonym / Restatement :} \eng{... leisure \textbf{or} free time...} / \eng{... hobbies, \textbf{in other words,} pastimes...}
\item \textbf{Antonym / Contrast :} \eng{Unlike work, \textbf{recreation} is relaxing.} (Marqueurs : \eng{unlike, however, but, whereas}).
\item \textbf{Example / Illustration :} \eng{Recreational activities \textbf{such as} football, dancing \textbf{and} reading...} / \eng{... activities \textbf{like} hiking and fishing.}
\item \textbf{General knowledge / Logic :} \eng{The farmer rested under a mango tree.} (Un manguier donne de l'ombre $\rightarrow$ \eng{rested} = se reposer).
\end{itemize}
\item \textbf{Substitute your guess.} Remplacez le mot inconnu par votre hypothèse (en français ou en anglais simple) et vérifiez si la phrase a du sens.
\item \textbf{Check word formation.} Préfixes (\eng{un-, re-, over-}), suffixes (\eng{-able, -tion, -ment}), racines connues.
\end{enumerate}
\end{methode}

\begin{exemplebox}[Practice guessing]
Phrase : \eng{The \textbf{choir} rehearsed every Saturday for the cultural day.}
\begin{enumerate}
\item \textbf{Part of speech :} Nom (sujet du verbe \eng{rehearsed}, article \eng{the} devant).
\item \textbf{Indices :} \eng{rehearsed} (répétait), \eng{cultural day} (jour culturel), \eng{every Saturday} (régularité).
\item \textbf{Hypothèse :} Groupe qui chante / chorale.
\item \textbf{Vérification :} \eng{The \textbf{choir} (chorale) rehearsed...} $\rightarrow$ Sens parfait.
\end{enumerate}
\end{exemplebox}

\begin{aretenir}[Les trois stratégies de lecture du Bac]
\begin{itemize}
\item \cle{Predicting} (avant la lecture) : anticiper le contenu grâce au titre, aux images et à ses connaissances.
\item \cle{Skimming} (lecture rapide) : lire vite pour saisir l'\textbf{idée générale} du texte — on répond à la question \eng{What is the text about?}
\item \cle{Scanning} (lecture ciblée) : chercher une \textbf{information précise} (un chiffre, un nom, une date) sans tout lire — on répond à la question \eng{Where does the text mention...?}
\item \cle{Guessing meaning} (inférence lexicale) : déduire le sens d'un mot inconnu grâce aux indices du contexte (définition, synonyme, antonyme, exemple, logique).
\end{itemize}
Nous les mettrons en pratique dans la section suivante avec le texte \eng{How We Spend Our Free Time}.
\end{aretenir}

\begin{exercice}[Your turn : predict, classify and guess]
\begin{enumerate}
\item Look at the title \eng{How We Spend Our Free Time} and write \textbf{two} predictions in English, beginning with \eng{The text will probably...} or \eng{It may...}.
\item Classify these activities as \emph{active} or \emph{passive}: \eng{dancing, watching films, gardening, listening to music, playing draughts}.
\item Complete with \eng{leisure}, \eng{hobby} or \eng{pastime}: \eng{My uncle's favourite \dots\ is breeding pigeons; he spends all his \dots\ time on it. For him, it is more than a \dots\ — it is a passion.}
\item \textbf{Guessing practice.} In the sentence \eng{“Unlike passive TV watching, hiking is a \textbf{strenuous} outdoor activity.”}, guess the meaning of \eng{strenuous}. Tick the correct synonym: \eng{a) easy \quad b) relaxing \quad c) demanding effort \quad d) dangerous}. Justify using a context clue.
\end{enumerate}
\end{exercice}

\begin{corrige}[Exercice “Your turn”]
\begin{enumerate}
\item Réponses libres, par exemple : \eng{The text will probably describe the free-time activities of young people.} / \eng{It may compare recreation in Cameroon and in other countries.}
\item \emph{Active} : \eng{dancing, gardening, playing draughts} (effort physique ou mental). \emph{Passive} : \eng{watching films, listening to music}.
\item \eng{My uncle's favourite \textbf{hobby} is breeding pigeons; he spends all his \textbf{leisure} time on it. For him, it is more than a \textbf{pastime} — it is a passion.} Rappel : \eng{hobby} désigne une activité régulière et passionnée ; \eng{leisure time} est la collocation correcte pour « temps libre » ; \eng{pastime} convient pour une simple occupation agréable.
\item \textbf{c) demanding effort.} Justification : l'indice de contraste \eng{Unlike passive TV watching} (contrairement à la télévision passive) signale que \eng{hiking} a une qualité opposée à \eng{passive} ; \eng{strenuous} signifie donc « qui demande de l'effort, fatigant ».
\end{enumerate}
\end{corrige}

Votre esprit est maintenant prêt : vous connaissez le vocabulaire du thème, vous savez classer les activités, vous maîtrisez les trois stratégies de lecture (predicting, skimming, scanning) et l'inférence lexicale. Passons à la lecture du texte \eng{How We Spend Our Free Time}.

\coursec{Reading text: “How We Spend Our Free Time”}

\courssub{From warm-up to reading: setting the scene}

Au cours de la section précédente, nous avons parlé de nos loisirs (\eng{free time activities}) : le football au quartier, la danse, la musique, les jeux de société, la télévision. Nous avons appris à poser la question \eng{What do you do in your free time?} et à répondre avec \eng{I like...}, \eng{I enjoy...}, \eng{I am fond of...}.

Nous allons maintenant \cle{lire un article} qui reprend exactement ce thème, mais à une échelle plus large : celui-ci compare les loisirs des jeunes dans \textbf{trois pays anglophones} : le Cameroun (la ville de Bamenda), le Nigeria (la ville de Lagos) et le Royaume-Uni. Avant de lire, posons-nous trois questions de prédiction (\emph{pre-reading questions}) :

\begin{enumerate}
\item \eng{What do you think teenagers in Lagos do after school?}
\item \eng{Do British teenagers have the same hobbies as you?}
\item \eng{Can recreation bring different cultures together?}
\end{enumerate}

Gardez vos réponses en tête : nous vérifierons après la lecture si vos prédictions étaient correctes. C'est la première stratégie du bon lecteur : \cle{anticiper} le contenu du texte avant de le lire.

\courssub{The reading text}

Lisez le texte une première fois \textbf{rapidement}, sans dictionnaire, pour en saisir l'idée générale (c'est la stratégie du \emph{skimming}, que nous étudierons en détail dans la section suivante).

\begin{textebox}[How We Spend Our Free Time — a magazine article]
In Bamenda, weekends are never dull. As soon as classes end on Friday, students rush to the football field behind the school. Some play in organised teams; others simply kick a ball around with their friends. On Saturday evenings, many young people rehearse traditional dances for the annual cultural festival, while families gather in the sitting room to play board games such as ludo and draughts. For these teenagers, leisure is above all a moment shared with the community.

Meanwhile in Lagos, many teenagers are passionate about music. Afrobeats songs play in every taxi, every shop and every home. After school, groups of friends meet to record their own songs on cheap smartphones, hoping to become the next big star. On Saturdays, the cinemas are full: Nollywood films attract huge crowds of young fans. Those who prefer to stay at home spend hours playing video games online with friends from other neighbourhoods.

In Britain, most teenagers belong to a sports club. They train twice a week and play matches at the weekend, whatever the weather. When they are not on the pitch, they are often indoors: streaming series and films has become the favourite leisure activity of British youth. Surprisingly, reading is still popular too; many adolescents belong to book clubs at school or at the local library.

Whatever the country, recreation brings people together. A football match in Bamenda, a concert in Lagos or a book club in London may look very different, but they all serve the same purpose: they help young people relax, make friends and feel part of a community. Leisure, in the end, is a language that everybody speaks.
\end{textebox}

\begin{definition}[Glossaire — les mots nouveaux du texte]
\begin{center}
\begin{tabularx}{\linewidth}{|l|l|X|}
\hline
\rowcolor{popblueL} \textbf{Mot anglais} & \textbf{Nature} & \textbf{Équivalent français} \\
\hline
dull & adjectif & ennuyeux, morne \\
\hline
to rush & verbe & se précipiter \\
\hline
meanwhile & adverbe & pendant ce temps \\
\hline
to rehearse & verbe & répéter (danse, théâtre, musique) \\
\hline
board game & nom & jeu de société \\
\hline
to gather & verbe & se rassembler, se réunir \\
\hline
to belong to & verbe & appartenir à, être membre de \\
\hline
leisure & nom & loisir(s), temps libre \\
\hline
whatever & pronom relatif & quel que soit \\
\hline
purpose & nom & but, objectif \\
\hline
\end{tabularx}
\end{center}
\end{definition}

\begin{attention}[Faux ami : « recreation » n'est pas « récréation »]
Le mot anglais \eng{recreation} se prononce avec un \emph{e} long au début : « ré-kri-É-cheun ». Il signifie \textbf{loisirs, activités de détente} en général, et non la « récréation » scolaire (la pause entre les cours). Pour la récréation à l'école, l'anglais dit \eng{break} ou \eng{break time}.

\begin{center}
\begin{tabularx}{\linewidth}{|X|X|X|}
\hline
\rowcolor{popblueL} Français & English & Remarque \\
\hline
les loisirs, les activités de détente & recreation & \eng{Recreation brings people together.} \\
\hline
la récréation (pause à l'école) & break, break time & \eng{The pupils play during break.} \\
\hline
recréer (créer de nouveau) & to re-create & autre mot, autre prononciation \\
\hline
\end{tabularx}
\end{center}

Attention aussi à \eng{to rehearse} (« ré-pé-ter ») : il s'emploie pour répéter une danse, une pièce ou un concert, jamais pour « répéter » une leçon à un élève (\eng{to repeat}).
\end{attention}

\begin{savaistu}[Nollywood, géant mondial du cinéma]
Nollywood, l'industrie cinématographique du Nigeria, produit chaque année \textbf{plus de films que Hollywood} en nombre de titres (même si les budgets sont bien plus modestes). C'est l'un des premiers employeurs culturels d'Afrique et ses films sont regardés dans toute l'Afrique anglophone, y compris dans les régions du Nord-Ouest et du Sud-Ouest du Cameroun. Le mot lui-même est un \emph{mot-valise} (\emph{blend}) formé de \emph{Nigeria} et \emph{Hollywood}, sur le modèle de \emph{Bollywood} (Bombay + Hollywood) pour le cinéma indien. Regarder des films de Nollywood est donc bien plus qu'un loisir : c'est une véritable industrie et un lien culturel entre les pays anglophones d'Afrique de l'Ouest.
\end{savaistu}

\courssub{Understanding the text: structure and key information}

Ce texte est un \cle{article comparatif} : sa structure est très régulière, ce qui facilite la lecture. Observons-la ensemble.

\begin{methode}[Repérer la structure d'un texte comparatif]
\begin{enumerate}
\item \textbf{Lire le titre} : il annonce le thème (\eng{How We Spend Our Free Time} = comment nous occupons notre temps libre).
\item \textbf{Repérer les mots-repères} (\emph{markers}) qui ouvrent chaque paragraphe : \eng{In Bamenda...}, \eng{Meanwhile in Lagos...}, \eng{In Britain...}, \eng{Whatever the country...}. Chaque repère annonce un changement de lieu ou d'étape.
\item \textbf{Associer chaque paragraphe à une idée} : paragraphe 1 = Bamenda ; paragraphe 2 = Lagos ; paragraphe 3 = la Grande-Bretagne ; paragraphe 4 = conclusion générale.
\item \textbf{Dégager la thèse de l'auteur} dans la conclusion : malgré les différences, les loisirs unissent les communautés (\eng{recreation brings people together}).
\end{enumerate}
\end{methode}

Appliquons cette méthode et résumons le contenu du texte dans un tableau comparatif.

\begin{center}
\begin{tabularx}{\linewidth}{|l|X|X|}
\hline
\rowcolor{popblueL} \textbf{Country} & \textbf{Activities mentioned} & \textbf{Place (indoor / outdoor)} \\
\hline
Cameroon (Bamenda) & football, traditional dance rehearsals, board games (ludo, draughts) & outdoor (football field) and indoor (sitting room) \\
\hline
Nigeria (Lagos) & afrobeats music, recording songs, Nollywood films, video games & indoor (cinemas, home) and outdoor (streets, taxis) \\
\hline
Britain & sports clubs, streaming series, reading, book clubs & outdoor (the pitch) and indoor (home, library) \\
\hline
\end{tabularx}
\end{center}

Ce tableau est déjà une forme de \cle{résumé visuel} du texte : il sera très utile pour répondre aux questions de compréhension et pour rédiger votre propre résumé dans la section 6.

\begin{popfigure}
\begin{center}\tcbox[colback=popblue,colframe=popblue,coltext=white,arc=9pt,boxsep=4pt,left=6pt,right=6pt]{\bfseries\small Recreation unites us}\end{center}
\vspace{-2pt}
\begin{tcbraster}[raster columns=2,raster equal height=rows,raster column skip=6pt,raster row skip=6pt,enhanced,blankest]
\begin{tcolorbox}[enhanced,colback=white,colframe=poporange,boxrule=0.9pt,arc=3pt,title={Bamenda},fonttitle=\bfseries\scriptsize,coltitle=white,colbacktitle=poporange,left=4pt,right=4pt,top=2pt,bottom=2pt,boxsep=1pt,toptitle=1.5pt,bottomtitle=1.5pt,halign=left]
\begin{itemize}[nosep,leftmargin=*,itemsep=1pt,font=\scriptsize]
\item football
\item ludo, draughts
\item traditional dance
\end{itemize}
\end{tcolorbox}
\begin{tcolorbox}[enhanced,colback=white,colframe=popgreen,boxrule=0.9pt,arc=3pt,title={Lagos},fonttitle=\bfseries\scriptsize,coltitle=white,colbacktitle=popgreen,left=4pt,right=4pt,top=2pt,bottom=2pt,boxsep=1pt,toptitle=1.5pt,bottomtitle=1.5pt,halign=left]
\begin{itemize}[nosep,leftmargin=*,itemsep=1pt,font=\scriptsize]
\item afrobeats
\item Nollywood
\item video games
\end{itemize}
\end{tcolorbox}
\begin{tcolorbox}[enhanced,colback=white,colframe=poppurple,boxrule=0.9pt,arc=3pt,title={Britain},fonttitle=\bfseries\scriptsize,coltitle=white,colbacktitle=poppurple,left=4pt,right=4pt,top=2pt,bottom=2pt,boxsep=1pt,toptitle=1.5pt,bottomtitle=1.5pt,halign=left]
\begin{itemize}[nosep,leftmargin=*,itemsep=1pt,font=\scriptsize]
\item sports clubs
\item streaming
\item book clubs
\end{itemize}
\end{tcolorbox}
\end{tcbraster}

\legende{Carte mentale du texte : trois pays, trois univers de loisirs, une même conclusion.}
\end{popfigure}

\courssub{First comprehension questions}

Répondez mentalement ou à l'écrit, en phrases complètes, à ces premières questions de compréhension globale (les questions détaillées et les stratégies de lecture feront l'objet des sections 3 et 4).

\begin{enumerate}
\item \eng{What is the main idea of the text?}
\item \eng{Which activity is shared by the three countries?}
\item \eng{According to the writer, why is recreation important?}
\item \eng{Which paragraph describes your own free time best? Why?}
\end{enumerate}

\begin{exoresolu}[Comprehension question — modèle de réponse]
Énoncé : \eng{According to the writer, why is recreation important? Answer in one or two complete sentences.}
\tcblower
\textbf{Méthode} : on repère d'abord le mot-clé de la question (\eng{important}) dans la conclusion du texte : \eng{they help young people relax, make friends and feel part of a community}. On reformule ensuite avec ses propres mots, sans recopier tout le paragraphe.

\textbf{Réponse modèle} : \eng{According to the writer, recreation is important because it helps young people relax, make friends and feel part of a community, whatever their country.}

« Selon l'auteur, les loisirs sont importants parce qu'ils aident les jeunes à se détendre, à se faire des amis et à se sentir membres d'une communauté, quel que soit leur pays. »

Remarque : la réponse commence par \eng{According to the writer...} pour montrer clairement qu'on rapporte l'idée du texte, et non sa propre opinion.
\end{exoresolu}

\begin{exercice}[Your turn — quick check]
Répondez en anglais, en phrases complètes :
\begin{enumerate}
\item \eng{Where do students in Bamenda play football?}
\item \eng{What do teenagers in Lagos hope to become when they record songs?}
\item \eng{How often do British teenagers train in their sports clubs?}
\item Trouvez dans le texte un synonyme de \eng{free time} et un contraire de \eng{exciting}.
\end{enumerate}
\end{exercice}

\begin{corrige}[Exercice « quick check »]
\begin{enumerate}
\item \eng{They play football on the football field behind the school.} (Ils jouent sur le terrain de football derrière l'école.)
\item \eng{They hope to become the next big star.} (Ils espèrent devenir la prochaine grande star.)
\item \eng{They train twice a week.} (Ils s'entraînent deux fois par semaine.)
\item Synonyme de \eng{free time} : \eng{leisure}. Contraire de \eng{exciting} : \eng{dull}.
\end{enumerate}
\end{corrige}

\begin{aretenir}[L'essentiel de cette section]
\begin{itemize}
\item Le texte \eng{How We Spend Our Free Time} est un \textbf{article comparatif} en quatre paragraphes : Bamenda, Lagos, Grande-Bretagne, conclusion.
\item Les \textbf{mots-repères} en début de paragraphe (\eng{In Bamenda...}, \eng{Meanwhile...}, \eng{Whatever the country...}) guident la lecture.
\item La \textbf{thèse de l'auteur} : les loisirs unissent les jeunes malgré les différences culturelles.
\item Vocabulaire clé : \eng{dull}, \eng{to rush}, \eng{meanwhile}, \eng{to rehearse}, \eng{board game}, \eng{to gather}, \eng{to belong to}, \eng{leisure}.
\item Piège : \eng{recreation} = « loisirs » (« ré-kri-É-cheun »), et non « récréation » scolaire (\eng{break}).
\end{itemize}
\end{aretenir}

Dans la section suivante, nous apprendrons à exploiter ce texte comme un lecteur expert grâce à trois stratégies : le \emph{skimming} (lecture rapide pour l'idée générale), le \emph{scanning} (recherche d'informations précises) et la \emph{déduction du sens des mots} à partir du contexte.

\coursec{Reading strategies: skimming, scanning, guessing}

Dans la section précédente, nous avons lu le texte \emph{“How We Spend Our Free Time”} et découvert comment des jeunes de Lagos, Yaoundé et Londres occupent leurs loisirs. Mais une question se pose : \textbf{comment} avons-nous lu ce texte ? Avons-nous traduit chaque mot ? Non — et c'est précisément ce que nous allons apprendre ici. Au baccalauréat, le candidat dispose d'un temps limité face à un texte inconnu. Trois stratégies de lecture professionnelles lui permettent de gagner du temps et des points : le \cle{skimming} (lecture globale), le \cle{scanning} (recherche ciblée) et la \cle{guessing} (déduction du sens des mots par le contexte).

Pour travailler ces stratégies, reprenons un extrait du texte support.

\begin{textebox}[How We Spend Our Free Time (extract)]
Weekends are never dull in Lagos. As soon as the last bell rings on Friday afternoon, students rush to the football field behind the school. Some play, others cheer, and a few simply sit under the mango tree to chat with friends. In Yaoundé, the picture is different but just as lively: many teenagers spend their Saturday mornings helping their parents at the market, and their afternoons watching football matches on television or playing board games such as ludo and draughts. In London, free time often means indoor activities because of the weather: visiting museums, which are usually free, going to the cinema, or meeting friends in a café. Yet, whatever the city, one thing is common everywhere: young people use their free time to relax, to socialise and to discover new interests. Recreation is not a waste of time; doctors say it recreates the mind and the body, and teachers agree that students who have enjoyable hobbies often perform better at school.
\end{textebox}

\textbf{Glossaire :} \emph{dull} (adj.) = ennuyeux ; \emph{to rush} (v.) = se précipiter ; \emph{to cheer} (v.) = encourager, applaudir ; \emph{board games} (n.) = jeux de société ; \emph{draughts} (n.) = jeu de dames ; \emph{whatever} (conj.) = quel que soit ; \emph{to perform} (v.) = ici, obtenir de bons résultats.

\courssub{Skimming: reading for the general idea}

\begin{definition}[Skimming]
Le \cle{skimming} (de l'anglais \eng{to skim}, « écrémer, survoler ») consiste à lire un texte \textbf{rapidement et globalement} pour en saisir l'idée générale (\eng{the gist}), sans s'arrêter sur les détails ni sur les mots inconnus. On « écrème » le texte : on ne prend que l'essentiel.
\end{definition}

Observez : si vous lisez seulement le \textbf{titre} (\emph{How We Spend Our Free Time}), la \textbf{première phrase} du texte (\eng{Weekends are never dull in Lagos.}) et la \textbf{dernière phrase} (\eng{...students who have enjoyable hobbies often perform better at school.}), vous savez déjà que le texte parle des \textbf{loisirs des jeunes dans différentes villes et de leur utilité}. Vous avez fait du skimming — en moins d'une minute !

\begin{methode}[Comment faire du skimming]
\begin{enumerate}
\item Lire le \textbf{titre} du texte : il annonce le thème.
\item Lire la \textbf{première phrase de chaque paragraphe} (la \eng{topic sentence}) : elle contient l'idée principale du paragraphe.
\item Lire la \textbf{dernière phrase du texte} : elle contient souvent la conclusion ou le message de l'auteur.
\item \textbf{Ne jamais s'arrêter} sur un mot inconnu, ne jamais ouvrir le dictionnaire à cette étape.
\item Formuler l'idée générale en \textbf{une seule phrase} : \eng{This text is about...}
\end{enumerate}
\end{methode}

\begin{exemplebox}[Skimming appliqué au texte support]
Question : \eng{What is the text about?}\par
Réponse en une phrase : \eng{The text is about how young people in Lagos, Yaoundé and London spend their free time, and why recreation is useful.}\par
Traduction : « Le texte parle de la façon dont les jeunes de Lagos, Yaoundé et Londres occupent leur temps libre, et de l'utilité des loisirs. »
\end{exemplebox}

\courssub{Scanning: looking for specific information}

\begin{definition}[Scanning]
Le \cle{scanning} consiste à chercher dans un texte une \textbf{information précise} — un nom propre, un lieu, une date, un chiffre — \textbf{sans tout lire}. L'œil parcourt la page rapidement, comme un scanner, et s'arrête uniquement sur le mot-clé recherché.
\end{definition}

\begin{methode}[Comment faire du scanning]
\begin{enumerate}
\item Lire la \textbf{question} et y repérer le \textbf{mot-clé} : un nom propre (\eng{Lagos}, \eng{Yaoundé}), un chiffre, une date, un lieu.
\item Chercher ce mot-clé \textbf{visuellement} dans le texte, ligne par ligne, sans lire les phrases entières.
\item Une fois le mot-clé trouvé, lire \textbf{la phrase qui l'entoure} (avant et après) : la réponse s'y trouve presque toujours.
\item Recopier ou reformuler la réponse en une phrase complète.
\end{enumerate}
\end{methode}

\begin{exoresolu}[Scanning : une question ciblée]
Énoncé : \eng{Which activity is popular in Lagos?} (Quelle activité est populaire à Lagos ?)\par
\tcblower
\textbf{Étape 1 — Mot-clé :} le nom propre \eng{Lagos}.\par
\textbf{Étape 2 — Recherche visuelle :} le mot \eng{Lagos} apparaît dès la première phrase du texte.\par
\textbf{Étape 3 — Lecture de la phrase voisine :} \eng{...students rush to the football field behind the school.}\par
\textbf{Étape 4 — Réponse en phrase complète :} \eng{In Lagos, playing football is a popular activity: students rush to the football field after school.}\par
Traduction : « À Lagos, le football est une activité populaire : les élèves se précipitent au terrain de football après les cours. »
\end{exoresolu}

\courssub{Guessing: deducing the meaning of unknown words}

Même après le skimming et le scanning, il restera toujours des mots inconnus. Faut-il paniquer ? Non. Un bon lecteur \textbf{déduit} le sens des mots grâce aux indices du contexte.

\begin{propriete}[Les quatre indices pour deviner le sens d'un mot]
Le sens d'un mot inconnu se déduit par :
\begin{enumerate}
\item \textbf{Le contexte} (les mots voisins) : synonymes, exemples, contrastes, définitions glissés dans la phrase. Exemple : \eng{Weekends are never \textbf{dull}... students \textbf{rush} to the football field.} Le contraste est clair : si les élèves \emph{se précipitent} joyeusement au terrain, les week-ends ne peuvent pas être « ennuyeux »... si, justement : \eng{dull} signifie « ennuyeux, morne », et la phrase dit qu'ils ne le sont \emph{jamais}.
\item \textbf{La famille du mot} : \eng{recreation} appartient à la famille de \eng{to recreate} (« recréer, ressourcer ») — le texte le confirme : \eng{it recreates the mind and the body}.
\item \textbf{Les préfixes et suffixes} : \eng{enjoy} + suffixe \eng{-able} (« qui peut être... ») = \eng{enjoyable}, « agréable, dont on peut jouir ». De même : \eng{use} + \eng{-ful} = \eng{useful} (utile) ; \eng{use} + \eng{-less} = \eng{useless} (inutile).
\item \textbf{La nature grammaticale du mot} : sa place dans la phrase indique s'il s'agit d'un nom, d'un verbe ou d'un adjectif, ce qui réduit les hypothèses.
\end{enumerate}
\end{propriete}

\begin{exemplebox}[Exemples traduits]
\eng{Weekends are never dull.}\par
« Les week-ends ne sont jamais ennuyeux. »\par
\eng{Students rush to the field.}\par
« Les élèves se précipitent au terrain. »\par
\eng{Recreation recreates the mind and the body.}\par
« Les loisirs ressourcent l'esprit et le corps. »\par
\eng{Reading is an enjoyable hobby.}\par
« La lecture est un passe-temps agréable. »
\end{exemplebox}

\begin{exoresolu}[Guessing : deviner un mot par le contexte]
Énoncé : \eng{Guess the meaning of the word “lively” in: “In Yaoundé, the picture is different but just as lively.”}\par
\tcblower
\textbf{Étape 1 — Nature du mot :} \eng{lively} est un adjectif (placé après \eng{is... as}).\par
\textbf{Étape 2 — Contexte :} la phrase suivante décrit des adolescents qui aident leurs parents au marché, regardent des matchs, jouent aux jeux de société : beaucoup d'activité, de mouvement.\par
\textbf{Étape 3 — Famille du mot :} \eng{lively} vient de \eng{life} (« vie ») + suffixe \eng{-ly}.\par
\textbf{Conclusion :} \eng{lively} = « animé, plein de vie ». Traduction de la phrase : « À Yaoundé, le tableau est différent mais tout aussi animé. »
\end{exoresolu}

\courssub{The winning order at the Bac}

Au baccalauréat, l'ordre des opérations est stratégique. On ne commence \textbf{jamais} par lire le texte de bout en bout : on commence par les \textbf{questions}, qui nous disent quoi chercher.

\begin{popfigure}
\begin{tikzpicture}[node distance=0.55cm, every node/.style={font=\small}]
\node[draw=popblue, fill=popblueL, rounded corners, align=center, text width=3.6cm] (q) {\textbf{1. Read the questions}\\ Repérer les mots-clés};
\node[draw=popgreen, fill=popgreenL, rounded corners, align=center, text width=3.6cm, below=of q] (sk) {\textbf{2. Skim for gist}\\ Titre + 1\textsuperscript{re} phrase + fin};
\node[draw=poporange, fill=poporangeL, rounded corners, align=center, text width=3.6cm, below=of sk] (sc) {\textbf{3. Scan for details}\\ Chercher les mots-clés};
\node[draw=poppurple, fill=poppurpleL, rounded corners, align=center, text width=3.6cm, below=of sc] (gu) {\textbf{4. Guess unknown words}\\ Contexte, famille, suffixes};
\node[draw=poppink, fill=poppinkL, rounded corners, align=center, text width=3.6cm, below=of gu] (an) {\textbf{5. Answer}\\ Phrases complètes};
\draw[-{Stealth}, thick, popdark] (q) -- (sk);
\draw[-{Stealth}, thick, popdark] (sk) -- (sc);
\draw[-{Stealth}, thick, popdark] (sc) -- (gu);
\draw[-{Stealth}, thick, popdark] (gu) -- (an);
\end{tikzpicture}
\legende{L'ordre conseillé au Bac : questions, skimming, scanning, guessing, réponse.}
\end{popfigure}

\begin{center}
\begin{tabularx}{\linewidth}{|X|X|X|}
\hline
\rowcolor{popblueL} Stratégie & Objectif & Ce qu'on lit réellement \\
\hline
Skimming & Idée générale (\eng{gist}) & Titre, premières phrases, conclusion \\
\hline
Scanning & Information précise & Uniquement les phrases autour du mot-clé \\
\hline
Guessing & Sens d'un mot inconnu & La phrase et les mots voisins \\
\hline
\end{tabularx}
\end{center}

\begin{attention}[L'erreur fatale : tout traduire mot à mot]
Ne traduisez \textbf{jamais} tout le texte mot à mot au baccalauréat : c'est une perte de temps énorme et cela ne rapporte aucun point. Les questions portent sur des informations précises, pas sur une traduction complète. Lisez \textbf{avec un objectif} : les questions d'abord, le texte ensuite. Autre piège : ne bloquez pas sur un mot inconnu — dans 90 \% des cas, il n'est pas nécessaire pour répondre, et dans les autres cas, le contexte le trahit.
\end{attention}

\begin{aretenir}[L'essentiel]
\begin{itemize}
\item \cle{Skimming} = lire vite pour l'idée générale (titre, premières phrases, dernière phrase).
\item \cle{Scanning} = chercher un mot-clé précis sans tout lire.
\item \cle{Guessing} = déduire le sens d'un mot par le contexte, la famille du mot et les préfixes/suffixes.
\item Ordre au Bac : \textbf{questions} $\rightarrow$ \textbf{skimming} $\rightarrow$ \textbf{scanning} $\rightarrow$ \textbf{guessing} $\rightarrow$ \textbf{réponse}.
\end{itemize}
\end{aretenir}

\begin{exercice}[À vous de jouer]
1. Skimming : en une phrase, dites quel est le sujet du texte support.\par
2. Scanning : \eng{Why do people in London often choose indoor activities?}\par
3. Guessing : devinez le sens de \eng{to socialise} dans la phrase \eng{young people use their free time to relax, to socialise and to discover new interests}.
\end{exercice}

\begin{corrige}[Exercice]
1. \eng{The text is about the recreational activities of young people in three cities and the benefits of free time.}\par
2. \eng{Because of the weather.} — mot-clé \eng{London} ; la phrase dit : \eng{indoor activities because of the weather} (« à cause du temps qu'il fait »).\par
3. \eng{To socialise} = « voir des gens, se lier avec les autres » (contexte : \eng{meeting friends}, \eng{chat with friends} ; famille : \eng{social} + suffixe verbal \eng{-ise}).
\end{corrige}

\coursec{Comprehension practice}

Dans la section précédente, nous avons appris à \emph{skim} (lire rapidement pour saisir l'idée générale), à \emph{scan} (repérer une information précise) et à deviner le sens des mots inconnus grâce au contexte. Nous allons maintenant mettre ces stratégies en action : c'est le moment de répondre à des \cle{questions de compréhension} sur le texte, exactement comme au baccalauréat. Chaque type de question exige une stratégie différente : c'est ce que nous allons découvrir, pas à pas.

\courssub{The reference text: a quick reminder}

Pour travailler efficacement, relisons le texte support de la leçon. Gardez-le sous les yeux pendant toute la séance : toutes les réponses doivent s'appuyer sur lui.

\begin{textebox}[How We Spend Our Free Time]
What do people do when they are not working or studying? In Britain, as in Cameroon, free time is precious, and people fill it in very different ways.

According to a recent survey of British teenagers, the most popular leisure activity is listening to music: almost nine out of ten teenagers say they listen to music every day, usually on their phones. Watching videos and using social media come second, while reading books is far less popular than it was twenty years ago. Surprisingly, only about half of the teenagers interviewed said they regularly play video games, although many people believe that all young people are gamers.

Sport remains an important part of recreation. Football is the favourite sport to watch, but swimming and cycling are the sports people actually practise most. Older people often prefer quieter hobbies, such as gardening, walking or playing cards with friends.

In my opinion, the most interesting result of the survey is that recreation brings people together. Whether it is a football match, a church choir or a video game played online with friends, leisure is rarely a solitary activity. I believe this is why hobbies matter so much: they do not only fill our free time, they also create friendships and strengthen communities.

Of course, free time is not equally distributed. Students in examination classes have little time to relax, and many young people in big cities complain that leisure activities are too expensive. A cinema ticket or a gym membership can cost a lot. Nevertheless, the survey shows that the simplest activities — chatting with friends, listening to the radio, playing draughts under a tree — are still the most loved.
\end{textebox}

\textbf{Glossaire.}

\begin{center}
\begin{tabularx}{\linewidth}{|X|X|X|}\hline
\rowcolor{popblueL} \textbf{English} & \textbf{Nature} & \textbf{Français} \\ \hline
survey & noun & enquête, sondage \\ \hline
leisure & noun & loisirs \\ \hline
gamer & noun & joueur de jeux vidéo \\ \hline
to practise a sport & verb & pratiquer un sport \\ \hline
choir & noun & chœur, chorale \\ \hline
solitary & adjective & solitaire \\ \hline
to strengthen & verb & renforcer \\ \hline
membership & noun & adhésion, abonnement \\ \hline
draughts & noun & jeu de dames \\ \hline
\end{tabularx}
\end{center}

\begin{attention}[Orthographe britannique : practise / practice]
En anglais britannique, \eng{to practise} (avec \textbf{s}) est le \textbf{verbe} : \eng{People practise swimming.} ; \eng{practice} (avec \textbf{c}) est le \textbf{nom} : \eng{Practice makes perfect.} (« C'est en forgeant qu'on devient forgeron. »). En anglais américain, on écrit \eng{practice} dans les deux cas. Au Bac camerounais, l'orthographe britannique est la norme attendue. Autre piège : \eng{leisure} se prononce « \emph{lè-jeur} » et \eng{choir} se prononce « \emph{kouaï-eur} » — ne vous fiez jamais à l'orthographe pour deviner la prononciation.
\end{attention}

\courssub{True or False questions: justifying with the text}

Observons d'abord un exemple typique de question au Bac :

\eng{British teenagers all play video games. True or False? Justify your answer with a line from the text.}

Beaucoup d'élèves répondent \eng{True}, parce que l'affirmation semble vraie dans la vie réelle. Erreur ! Il faut répondre \textbf{d'après le texte}, pas d'après ses opinions personnelles. Le texte dit : \eng{only about half of the teenagers interviewed said they regularly play video games}. La réponse correcte est donc :

\eng{False. The text says: “only about half of the teenagers interviewed said they regularly play video games.”}

\begin{methode}[Répondre au Bac : les règles d'or]
\begin{enumerate}
\item \textbf{Lisez les questions avant le texte} si l'examen le permet : vous saurez quoi chercher (scanning).
\item \textbf{Justifiez toujours} les réponses \emph{True/False} en citant le passage exact du texte entre guillemets. Une réponse sans justification perd la moitié des points, même si elle est correcte.
\item \textbf{Répondez par des phrases complètes} avec sujet + verbe, sauf si la consigne demande explicitement un mot ou une expression (\eng{Answer in one word}).
\item \textbf{Ne recopiez jamais tout un paragraphe} : citez uniquement la ligne qui prouve votre réponse.
\item \textbf{Basez-vous sur le texte, pas sur vos connaissances} : si le texte contredit la réalité, c'est le texte qui a raison.
\end{enumerate}
\end{methode}

\begin{exoresolu}[True or False with justification]
\textbf{Énoncé.} Say whether the following statements are True or False. Justify each answer with a quotation from the text.

1. \eng{Reading books is more popular than twenty years ago.}

2. \eng{Football is the sport people practise most.}

3. \eng{Some leisure activities cost a lot of money.}
\tcblower
\textbf{Solution détaillée.}

1. \eng{False. The text says: “reading books is far less popular than it was twenty years ago.”} — Le texte affirme le contraire : la lecture a \emph{diminué}.

2. \eng{False. The text says: “Football is the favourite sport to watch, but swimming and cycling are the sports people actually practise most.”} — Piège classique : le football est le sport le plus \emph{regardé}, pas le plus \emph{pratiqué}.

3. \eng{True. The text says: “A cinema ticket or a gym membership can cost a lot.”} — La justification est directe et courte : une ligne suffit.
\end{exoresolu}

\courssub{Multiple choice questions (MCQ)}

Le QCM (\eng{multiple choice questions}) propose quatre options : une correcte et trois \cle{distracteurs}. Observons un exemple :

\eng{According to the survey, the most popular leisure activity among British teenagers is:}

\begin{enumerate}[label=\Alph*.]
\item \eng{playing video games}
\item \eng{listening to music}
\item \eng{reading books}
\item \eng{watching football}
\end{enumerate}

La bonne réponse est \textbf{B} : \eng{the most popular leisure activity is listening to music}. Remarquez que les options A, C et D contiennent toutes des mots qui \emph{apparaissent dans le texte} (\eng{video games}, \eng{reading books}, \eng{football}) : c'est le principe du distracteur.

\begin{attention}[Les distracteurs : des mots du texte, un sens déformé]
Dans les QCM, les mauvaises options reprennent presque toujours des mots du texte pour vous attirer, mais elles en \textbf{déforment le sens}. Ne choisissez jamais une option simplement parce que vous avez vu ses mots dans le texte : \textbf{relisez le contexte complet} de la phrase concernée et vérifiez que l'option dit exactement la même chose. Comparez : le texte dit \eng{Football is the favourite sport to watch} ; un distracteur dira \eng{Football is the most practised sport} — les mots sont là, le sens est faux.
\end{attention}

\begin{exoresolu}[MCQ: details and general idea]
\textbf{Énoncé.} Choose the correct option (A, B, C or D).

1. \eng{Older people often prefer:}
A. \eng{video games} \quad B. \eng{social media} \quad C. \eng{quieter hobbies like gardening} \quad D. \eng{swimming competitions}

2. \eng{The general idea of the text is:}
A. \eng{British teenagers are lazy} \quad B. \eng{how people use their free time and why it matters} \quad C. \eng{leisure activities are too expensive} \quad D. \eng{reading is dead}
\tcblower
\textbf{Solution détaillée.}

1. Réponse \textbf{C}. Le texte dit : \eng{Older people often prefer quieter hobbies, such as gardening, walking or playing cards}. Les options A et B décrivent les adolescents ; D n'est jamais mentionné.

2. Réponse \textbf{B}. C'est une question d'\emph{idée générale} : utilisez le skimming (titre + première et dernière phrases). L'option A est un jugement absent du texte ; C n'est qu'un détail du dernier paragraphe ; D est une exagération (\eng{far less popular} ne veut pas dire \eng{dead}).
\end{exoresolu}

\courssub{Wh-questions: short answers in complete sentences}

Les questions en \emph{wh-} (\eng{Who? What? Where? When? Why? How?}) exigent une réponse courte mais complète. Voici comment repérer ce que chaque question demande :

\begin{center}
\begin{tabularx}{\linewidth}{|X|X|X|}\hline
\rowcolor{popblueL} \textbf{Question word} & \textbf{Elle demande...} & \textbf{Exemple de réponse} \\ \hline
\eng{Who...?} & une personne, un groupe & \eng{British teenagers.} \\ \hline
\eng{What...?} & une chose, une activité & \eng{Listening to music.} \\ \hline
\eng{Where...?} & un lieu & \eng{In Britain.} \\ \hline
\eng{When...?} & un moment & \eng{Twenty years ago.} \\ \hline
\eng{Why...?} & une raison (\eng{because...}) & \eng{Because they create friendships.} \\ \hline
\eng{How...?} & une manière, un moyen & \eng{On their phones.} \\ \hline
\end{tabularx}
\end{center}

\begin{exemplebox}[Wh-questions sur le texte]
\eng{Who listens to music every day?} — \eng{Almost nine out of ten British teenagers listen to music every day.} (Presque neuf adolescents britanniques sur dix écoutent de la musique chaque jour.)

\eng{What sports do people actually practise most?} — \eng{They practise swimming and cycling most.} (Ils pratiquent surtout la natation et le cyclisme.)

\eng{Why do some young people complain?} — \eng{They complain because leisure activities are too expensive.} (Ils se plaignent parce que les loisirs coûtent trop cher.)
\end{exemplebox}

\begin{attention}[Répondez à la question posée, pas à côté]
À \eng{Why...?}, la réponse doit contenir \eng{because} : répondre \eng{Leisure activities are expensive} sans \eng{because} est toléré, mais la forme complète est plus sûre. À \eng{Who...?}, ne répondez pas par un lieu. Et évitez le calque du français : on dit \eng{They complain because activities are expensive}, jamais \eng{They complain because of activities are expensive} — \eng{because of} est suivi d'un \textbf{nom} (\eng{because of the price}), \eng{because} d'une \textbf{proposition complète} (\eng{because the price is high}).
\end{attention}

\courssub{Interpretation and inference: reading between the lines}

Certaines questions ne demandent pas de trouver une information, mais d'\textbf{interpréter} le texte. Observons :

\eng{What does the author mean by “recreation brings people together”?}

L'expression \eng{recreation brings people together} se traduit : « Les loisirs rassemblent les gens ». L'auteur explique lui-même sa pensée : les loisirs sont rarement solitaires — un match de football, une chorale, un jeu vidéo en ligne se vivent avec d'autres — et ils « créent des amitiés et renforcent les communautés ». Une bonne réponse serait :

\eng{The author means that leisure activities are usually shared with other people and that they help create friendships and strengthen communities.}

\begin{propriete}[Les questions d'inférence]
Une question d'inférence commence souvent par \eng{What can we deduce...?}, \eng{What does the author suggest...?}, \eng{What does the author mean by...?}. Elle exige de \textbf{lire entre les lignes} : la réponse n'est pas écrite mot pour mot dans le texte. Il faut combiner plusieurs indices et les reformuler avec ses propres mots. Exemple : \eng{What can we deduce about students in examination classes?} — Le texte dit qu'ils \eng{have little time to relax} ; on peut donc déduire : \eng{We can deduce that they are under pressure and cannot enjoy leisure as much as others.}
\end{propriete}

\courssub{Fact or opinion?}

Un lecteur compétent distingue ce qui est \textbf{vérifiable} (un \cle{fait}) de ce qui est un \textbf{jugement personnel} (une \cle{opinion}). Les opinions sont souvent signalées par \eng{I believe}, \eng{in my opinion}, \eng{the most interesting}, \eng{surprisingly}.

\begin{center}
\begin{tabularx}{\linewidth}{|X|X|X|}\hline
\rowcolor{popblueL} \textbf{Phrase du texte} & \textbf{Fait ou opinion ?} & \textbf{Indice} \\ \hline
\eng{Almost nine out of ten teenagers listen to music every day.} & Fait & chiffre vérifiable (enquête) \\ \hline
\eng{The most interesting result is that recreation brings people together.} & Opinion & \eng{in my opinion}, jugement de valeur \\ \hline
\eng{Swimming and cycling are the sports people practise most.} & Fait & résultat d'enquête \\ \hline
\eng{Hobbies matter so much.} & Opinion & \eng{I believe}, appréciation personnelle \\ \hline
\end{tabularx}
\end{center}

\begin{exemplebox}[Exemple traduit]
\eng{Recreation brings people together.} = « Les loisirs rassemblent les gens. »

C'est une \textbf{opinion de l'auteur}, pas un fait vérifiable : on ne peut pas la mesurer comme un pourcentage. L'auteur la présente avec \eng{In my opinion} et \eng{I believe}, ce qui signale clairement un point de vue personnel.
\end{exemplebox}

\courssub{Synonyms in context: reformulation}

Au Bac, on vous demande souvent de \textbf{trouver dans le texte un synonyme} d'un mot donné (\eng{Find in the text a word that means...}). Stratégie : repérez la zone du texte où le thème apparaît, puis cherchez un mot de même nature (nom, adjectif, adverbe) et de sens proche. Le mot doit \textbf{exister réellement dans le texte} : ne l'inventez jamais.

\begin{exoresolu}[Find synonyms in the text]
\textbf{Énoncé.} Find in the text words or expressions that mean the same as:

1. \eng{free time} \quad 2. \eng{alone} \quad 3. \eng{to make stronger}
\tcblower
\textbf{Solution détaillée.}

1. \eng{free time} (temps libre) : le texte utilise \eng{leisure} et \eng{recreation} — tous deux signifient « loisirs ». Réponse : \eng{leisure} (ou \eng{recreation}).

2. \eng{alone} (seul) : le texte contient l'adjectif \eng{solitary} dans \eng{leisure is rarely a solitary activity}. Réponse : \eng{solitary}.

3. \eng{to make stronger} (rendre plus fort) : le texte contient le verbe \eng{strengthen} dans \eng{they also create friendships and strengthen communities}. Réponse : \eng{strengthen}.

Leçon de méthode : les trois réponses sont des mots \textbf{réellement présents} dans le texte, de même nature grammaticale que le mot de la consigne (nom pour nom, adjectif pour adjectif, verbe pour verbe). Si vous ne trouvez pas le synonyme, relisez la zone thématique ligne par ligne avec le scanning — mais ne proposez jamais un mot absent du texte.
\end{exoresolu}

\courssub{Choosing the right strategy for each question}

Chaque type de question mobilise une stratégie de lecture précise. Voici le tableau de correspondance à mémoriser :

\begin{popfigure}
\begin{center}
\begin{tabularx}{\linewidth}{|X|X|X|}
\hline
\rowcolor{popblueL} \textbf{Type de question} & \textbf{Ce qu'elle teste} & \textbf{Stratégie à utiliser} \\
\hline
\eng{True/False + justification} & compréhension de détails précis & \textbf{scanning} : repérer la ligne exacte et la citer \\
\hline
QCM sur un détail & compréhension fine d'un passage & \textbf{scanning} + relecture du contexte (méfiez-vous des distracteurs) \\
\hline
QCM sur l'idée générale & compréhension globale & \textbf{skimming} : titre, introduction, conclusion \\
\hline
\eng{Wh-questions} & repérage d'informations factuelles & \textbf{scanning} guidé par le mot interrogatif \\
\hline
Inférence / interprétation & lecture entre les lignes & \textbf{skimming} + recoupement de plusieurs indices \\
\hline
Synonymes & richesse du vocabulaire en contexte & \textbf{scanning} dans la zone thématique + deviner par le contexte \\
\hline
Résumé à trous & compréhension globale + grammaire & \textbf{skimming} puis vérification grammaticale de chaque trou \\
\hline
\end{tabularx}
\end{center}
\legende{À chaque type de question, sa stratégie de lecture.}
\end{popfigure}

\courssub{Guided summary: completing a gapped summary}

Dernier exercice typique : le \cle{résumé à trous}. Vous devez compléter un résumé du texte avec des mots du texte ou de votre propre vocabulaire. Vérifiez toujours que le mot choisi est \textbf{grammaticalement correct} dans la phrase (nature du mot, accord, temps).

\begin{exoresolu}[Gapped summary in three sentences]
\textbf{Énoncé.} Complete the summary with words from the text.

\eng{The text is a ............ about how people in Britain spend their ............ time. A survey shows that teenagers prefer listening to ............, while older people enjoy quieter ............ such as gardening. The author believes that recreation brings people ............ because it creates friendships and strengthens ............ .}
\tcblower
\textbf{Solution détaillée.}

1. \eng{report} — nom singulier après l'article \eng{a} ; le texte est un compte rendu d'enquête (\eng{survey report}).

2. \eng{free} — l'expression du texte est \eng{free time} (« temps libre »).

3. \eng{music} — détail du deuxième paragraphe ; nom indénombrable, sans article ni pluriel.

4. \eng{hobbies} — pluriel exigé par \eng{such as gardening} et l'absence d'article.

5. \eng{together} — citation directe de la thèse de l'auteur (\eng{recreation brings people together}).

6. \eng{communities} — pluriel, comme dans le texte (\eng{strengthen communities}).

Résumé complet : \eng{The text is a report about how people in Britain spend their free time. A survey shows that teenagers prefer listening to music, while older people enjoy quieter hobbies such as gardening. The author believes that recreation brings people together because it creates friendships and strengthens communities.}
\end{exoresolu}

\begin{exercice}[À vous de jouer]
1. Say True or False and justify: \eng{Free time is equally distributed among all students.}

2. Choose the correct option: \eng{The simplest activities are:} A. \eng{the most expensive} \quad B. \eng{the most loved} \quad C. \eng{the least popular} \quad D. \eng{only for old people}

3. Answer in complete sentences: \eng{Why do many young people in big cities complain? What are the simplest activities mentioned in the text?}

4. Fact or opinion? \eng{“Leisure is rarely a solitary activity.”} Justify your choice.

5. Find in the text a word that means \eng{alone}.
\end{exercice}

\begin{corrige}[Exercice : À vous de jouer]
1. \eng{False. The text says: “free time is not equally distributed.”}

2. Réponse \textbf{B} : \eng{the simplest activities ... are still the most loved}. Les options A et C déforment le sens ; D n'est pas dit.

3. \eng{They complain because leisure activities are too expensive. The simplest activities mentioned are chatting with friends, listening to the radio and playing draughts under a tree.}

4. C'est une \textbf{opinion} (une généralisation de l'auteur), introduite dans un paragraphe marqué par \eng{In my opinion} ; elle n'est accompagnée d'aucun chiffre vérifiable.

5. \eng{solitary} (« solitaire », synonyme de \eng{alone}).
\end{corrige}

\begin{aretenir}[L'essentiel de la section]
\begin{itemize}
\item \textbf{True/False} : toujours citer la ligne du texte qui justifie la réponse.
\item \textbf{QCM} : les distracteurs reprennent les mots du texte avec un sens déformé — vérifier le contexte.
\item \textbf{Wh-questions} : réponses courtes mais en phrases complètes ; \eng{Why?} appelle \eng{because} + proposition, \eng{because of} + nom.
\item \textbf{Inférence} : lire entre les lignes, reformuler avec ses propres mots.
\item \textbf{Fait vs opinion} : repérer \eng{I believe}, \eng{in my opinion}, les jugements de valeur.
\item \textbf{Synonymes} : ne proposer que des mots réellement présents dans le texte, de même nature grammaticale.
\item \textbf{Résumé à trous} : le mot choisi doit être correct en sens \emph{et} en grammaire.
\end{itemize}
\end{aretenir}

\coursec{Vocabulary: the world of recreation}

Après avoir travaillé la compréhension du texte \eng{How We Spend Our Free Time}, nous allons maintenant enrichir notre vocabulaire. Comprendre un texte, c'est bien ; pouvoir parler et écrire sur le même sujet avec des mots précis et des constructions correctes, c'est l'objectif de l'examen. Cette section vous donne tout le lexique nécessaire pour décrire vos loisirs, ceux de votre famille et ceux des jeunes Camerounais, avec les bonnes structures grammaticales.

\courssub{Observing the words in context}

Avant toute liste de vocabulaire, observons les mots dans un texte authentique. Lisez ce court article et repérez les mots qui désignent des loisirs, ainsi que les verbes qui les accompagnent (\eng{play}, \eng{go}, \eng{do}, \eng{watch}, \eng{listen to}...).

\begin{textebox}[Weekend Life in Yaoundé — a school magazine article]
Ask any student at Lycée de Nkol-Eton how they spend the weekend, and you will hear the same answer: there is never enough free time! On Saturday mornings, many boys play football on the dusty pitch behind the school, while others prefer to stay at home and play video games or board games such as ludo and chess with their brothers and sisters.

Music is everywhere in the city. Some teenagers belong to the church choir and practise singing on Saturday afternoons. Others are fond of dancing and join drama clubs, where they rehearse plays for the end-of-year concert. Those who are keen on painting meet at the cultural centre to work on their pictures.

Not everyone likes active hobbies. Many students simply watch TV, stream films on their phones, or browse the internet in a cybercafé. Reading novels is less common, but a small group of passionate readers meets every month at the municipal library to discuss African literature.

On Sundays, families often choose outdoor activities. Some go swimming at the municipal pool, others go fishing on the river, and a few energetic ones go hiking on the hills around the city. Gardening is also popular: many parents spend hours planting flowers and vegetables. And of course, the traditional family picnic remains the favourite way to relax together before a new school week begins.

Whatever the activity, one thing is certain: recreational activities help young people to rest, to make friends and to discover their talents.
\end{textebox}

\textbf{Glossaire}

\begin{center}
\begin{tabularx}{\linewidth}{|l|l|X|}
\hline
\rowcolor{popblueL} Mot anglais & Nature & Traduction française \\
\hline
pitch & nom & terrain (de sport) \\
\hline
board games & nom pluriel & jeux de société \\
\hline
choir & nom & chœur, chorale \\
\hline
to rehearse & verbe & répéter (une pièce, un concert) \\
\hline
to stream & verbe & regarder en streaming \\
\hline
to browse the internet & expression & naviguer sur internet \\
\hline
cybercafé & nom & cybercafé \\
\hline
hiking & nom & randonnée \\
\hline
pool & nom & piscine \\
\hline
picnic & nom & pique-nique \\
\hline
talent & nom & don, talent \\
\hline
\end{tabularx}
\end{center}

\courssub{The four families of recreational activities}

Classons maintenant le vocabulaire du texte (et quelques mots utiles supplémentaires) en quatre grandes familles. Ce classement vous aidera à mémoriser les mots et à organiser vos productions écrites et orales.

\begin{center}
\begin{tabularx}{\linewidth}{|X|X|X|X|}
\hline
\rowcolor{popblueL} Sports \& Games & Arts \& Music & Media & Outdoor activities \\
\hline
football — le football & music — la musique & watching TV — regarder la télévision & hiking — la randonnée \\
\hline
athletics — l'athlétisme & dancing — la danse & streaming — le streaming & swimming — la natation \\
\hline
board games — les jeux de société & drama — le théâtre & reading novels — lire des romans & fishing — la pêche \\
\hline
chess — les échecs & concert — le concert & browsing the internet — naviguer sur internet & picnicking — le pique-nique \\
\hline
ludo — le jeu de ludo (petits chevaux) & choir — la chorale & listening to music — écouter de la musique & gardening — le jardinage \\
\hline
video games — les jeux vidéo & cinema — le cinéma & reading magazines — lire des magazines & camping — le camping \\
\hline
volleyball — le volley-ball & painting — la peinture & using a smartphone — utiliser son smartphone & cycling — le cyclisme \\
\hline
\end{tabularx}
\end{center}

\begin{popfigure}
\begin{center}\tcbox[colback=poporange!70,colframe=poporange!70,coltext=white,arc=9pt,boxsep=4pt,left=6pt,right=6pt]{\bfseries\small Recreation}\end{center}
\vspace{-2pt}
\begin{tcbraster}[raster columns=2,raster equal height=rows,raster column skip=6pt,raster row skip=6pt,enhanced,blankest]
\begin{tcolorbox}[enhanced,colback=white,colframe=popblue!50,boxrule=0.9pt,arc=3pt,title={Sports \& Games},fonttitle=\bfseries\scriptsize,coltitle=white,colbacktitle=popblue!50,left=4pt,right=4pt,top=2pt,bottom=2pt,boxsep=1pt,toptitle=1.5pt,bottomtitle=1.5pt,halign=left]
\begin{itemize}[nosep,leftmargin=*,itemsep=1pt,font=\scriptsize]
\item football
\item chess
\end{itemize}
\end{tcolorbox}
\begin{tcolorbox}[enhanced,colback=white,colframe=poppurple!50,boxrule=0.9pt,arc=3pt,title={Arts \& Music},fonttitle=\bfseries\scriptsize,coltitle=white,colbacktitle=poppurple!50,left=4pt,right=4pt,top=2pt,bottom=2pt,boxsep=1pt,toptitle=1.5pt,bottomtitle=1.5pt,halign=left]
\begin{itemize}[nosep,leftmargin=*,itemsep=1pt,font=\scriptsize]
\item choir
\item drama
\end{itemize}
\end{tcolorbox}
\begin{tcolorbox}[enhanced,colback=white,colframe=poppink!50,boxrule=0.9pt,arc=3pt,title={Media},fonttitle=\bfseries\scriptsize,coltitle=white,colbacktitle=poppink!50,left=4pt,right=4pt,top=2pt,bottom=2pt,boxsep=1pt,toptitle=1.5pt,bottomtitle=1.5pt,halign=left]
\begin{itemize}[nosep,leftmargin=*,itemsep=1pt,font=\scriptsize]
\item TV
\item internet
\end{itemize}
\end{tcolorbox}
\begin{tcolorbox}[enhanced,colback=white,colframe=popgreen!50,boxrule=0.9pt,arc=3pt,title={Outdoor},fonttitle=\bfseries\scriptsize,coltitle=white,colbacktitle=popgreen!50,left=4pt,right=4pt,top=2pt,bottom=2pt,boxsep=1pt,toptitle=1.5pt,bottomtitle=1.5pt,halign=left]
\begin{itemize}[nosep,leftmargin=*,itemsep=1pt,font=\scriptsize]
\item hiking
\item fishing
\end{itemize}
\end{tcolorbox}
\end{tcbraster}

\legende{Carte mentale : les quatre familles du monde des loisirs.}
\end{popfigure}

\courssub{The golden rule: PLAY, DO or GO?}

En français, on dit simplement « faire du foot », « faire de la natation », « faire du judo » : un seul verbe, \emph{faire}. L'anglais, lui, utilise \textbf{trois verbes différents} selon le type d'activité. C'est l'une des règles les plus testées à l'examen du Baccalauréat.

\begin{propriete}[PLAY, DO, GO : la règle des trois verbes]
\begin{itemize}
\item \cle{PLAY} + sport avec ballon/ballon ovale ou jeu (avec ou sans adversaire) : \eng{play football}, \eng{play volleyball}, \eng{play chess}, \eng{play ludo}, \eng{play video games}, \eng{play the piano} (pour les instruments de musique aussi !).
\item \cle{DO} + activité individuelle, sport de combat, art martial ou activité de forme/entraînement : \eng{do athletics}, \eng{do karate}, \eng{do judo}, \eng{do gymnastics}, \eng{do yoga}, \eng{do exercise}.
\item \cle{GO} + verbe en \textbf{-ing} pour les activités de déplacement, de plein air ou se terminant par -ing : \eng{go swimming}, \eng{go hiking}, \eng{go fishing}, \eng{go dancing}, \eng{go cycling}, \eng{go shopping}, \eng{go camping}.
\end{itemize}
Exceptions à noter : on dit \eng{watch a match} (regarder un match), \eng{attend a concert} (assister à un concert), \eng{listen to music} (écouter de la musique) — ces verbes ne suivent pas la règle PLAY/DO/GO.
\end{propriete}

\begin{center}
\begin{tabular}{|l|l|l|}
\hline
\rowcolor{popgreenL} PLAY & DO & GO + -ing \\
\hline
play football & do athletics & go swimming \\
\hline
play chess & do karate & go hiking \\
\hline
play ludo & do judo & go fishing \\
\hline
play video games & do gymnastics & go dancing \\
\hline
play the guitar & do yoga & go cycling \\
\hline
\end{tabular}
\end{center}

\begin{exemplebox}[Exemples traduits]
\eng{I play football every Saturday with my friends.} « Je joue au football tous les samedis avec mes amis. »

\eng{My sister does gymnastics at the stadium.} « Ma sœur fait de la gymnastique au stade. »

\eng{We go fishing on the Wouri river during the holidays.} « Nous allons pêcher sur le fleuve Wouri pendant les vacances. »

\eng{They played chess all evening.} « Ils ont joué aux échecs toute la soirée. »
\end{exemplebox}

\begin{attention}[Les pièges des francophones]
\begin{itemize}
\item On dit \eng{to watch TV} et jamais \eng{to look TV}. \emph{Look} exige la préposition \emph{at} : \eng{look at the picture} (regarde l'image).
\item On dit \eng{to listen \textbf{to} music} : la préposition \emph{to} est obligatoire, alors qu'en français on « écoute de la musique » sans préposition.
\item Ne dites pas \eng{I make sport} : le verbe français « faire » ne se traduit pas toujours par \emph{make}. Dites \eng{I play football} ou \eng{I do athletics}.
\item \eng{to practise} s'écrit avec un \emph{s} (verbe) en anglais britannique : \eng{The choir practises on Saturdays.} (Le nom s'écrit \eng{practice} avec un \emph{c}).
\end{itemize}
\end{attention}

\courssub{Expressing your tastes: useful expressions}

Pour parler de vos loisirs, il ne suffit pas de connaître les noms des activités : il faut aussi savoir dire que vous les aimez. Voici les expressions essentielles, classées de la plus faible à la plus forte intensité.

\begin{center}
\begin{tabularx}{\linewidth}{|X|X|X|}
\hline
\rowcolor{popblueL} Expression anglaise & Traduction française & Exemple \\
\hline
in my free time / in my spare time & pendant mon temps libre & In my free time, I read novels. \\
\hline
I like + -ing / to + infinitif & j'aime & I like dancing. / I like to dance. \\
\hline
I enjoy + -ing & j'aime (plaisir ressenti) & I enjoy reading in the evening. \\
\hline
to be fond of + -ing & aimer beaucoup & She is fond of gardening. \\
\hline
to be keen on + -ing & être passionné de, adorer & I am keen on dancing. \\
\hline
to be passionate about + nom/-ing & être passionné par & He is passionate about music. \\
\hline
to spend time + -ing & passer du temps à & She spends her free time reading. \\
\hline
\end{tabularx}
\end{center}

\begin{exemplebox}[Exemples traduits — à retenir par cœur]
\eng{I am keen on dancing.} « J'adore danser. »

\eng{She spends her free time listening to music.} « Elle passe son temps libre à écouter de la musique. »

\eng{We are fond of board games like ludo and chess.} « Nous aimons beaucoup les jeux de société comme le ludo et les échecs. »

\eng{My brother is passionate about video games.} « Mon frère est passionné de jeux vidéo. »
\end{exemplebox}

\begin{attention}[Construction obligatoire : la préposition + le verbe en -ing]
Après \eng{be fond of}, \eng{be keen on}, \eng{be passionate about} et après \textbf{toute préposition} en anglais, le verbe qui suit prend la terminaison \textbf{-ing} (gérondif). On dit \eng{I am keen on danc\textbf{ing}} et jamais \eng{I am keen on to dance}. De même : \eng{I am interested in swimm\textbf{ing}}, \eng{I am good at play\textbf{ing} chess}, \eng{I look forward to see\textbf{ing} you}.
\end{attention}

\courssub{Word families: building more words from one}

Les examinateurs adorent tester les familles de mots (dérivation). À partir d'un seul radical, vous pouvez former un nom, un verbe, un adjectif. Voici les trois familles les plus utiles pour ce chapitre.

\begin{center}
\begin{tabularx}{\linewidth}{|l|l|X|}
\hline
\rowcolor{popgoldL} Nom & Adjectif / Verbe & Exemple traduit \\
\hline
recreation (les loisirs) & recreational (adj., de loisirs) & \eng{Football is a recreational activity.} « Le football est une activité de loisirs. » \\
\hline
enjoyment (le plaisir) & enjoyable (adj., agréable) ; to enjoy (v., aimer) & \eng{The picnic was very enjoyable.} « Le pique-nique était très agréable. » \\
\hline
relaxation (la détente) & relaxing (adj., relaxant) ; to relax (v., se détendre) & \eng{Gardening is relaxing.} « Le jardinage est relaxant. » \\
\hline
\end{tabularx}
\end{center}

\begin{attention}[Faux ami célèbre : library ou bookshop ?]
\eng{A library} est une \textbf{bibliothèque} : l'endroit où l'on \emph{emprunte} des livres gratuitement. Une \textbf{librairie}, l'endroit où l'on \emph{achète} des livres, se dit \eng{a bookshop} (UK) ou \eng{a bookstore} (US). Ne dites donc jamais \eng{I bought this novel at the library} : cela voudrait dire que vous avez acheté un livre à la bibliothèque, ce qui est impossible !
\end{attention}

\begin{savaistu}[How do young English speakers say “se détendre” ?]
L'expression \eng{to chill out} (prononcez « tchile aoute ») est extrêmement courante chez les jeunes anglophones pour dire « se détendre, se relaxer, ne rien faire de spécial ». \eng{I'm just chilling out at home.} « Je me détends simplement à la maison. » On entend aussi \eng{to hang out with friends} (traîner, passer du temps avec des amis, sans activité précise). Attention : ces expressions sont familières (\emph{informal}) — parfaites à l'oral entre amis, mais à éviter dans une rédaction formelle d'examen (composition, lettre formelle) !
\end{savaistu}

\begin{exoresolu}[Exercice : PLAY, DO ou GO ?]
Complétez chaque phrase avec \eng{play}, \eng{do} ou \eng{go} à la forme correcte (conjugaison au présent).

1. Every Sunday, my father \ldots\ fishing on the river Nyong.

2. The students of our school \ldots\ athletics twice a week.

3. My friends and I \ldots\ ludo during the break.

4. In her free time, Amina \ldots\ gymnastics at the stadium.

5. We often \ldots\ swimming at the municipal pool in Douala.

6. My uncle can \ldots\ the guitar very well.
\tcblower
\textbf{Solution détaillée}

1. \eng{goes} — \eng{fishing} est un verbe en -ing, donc GO ; sujet \eng{my father} (3\textsuperscript{e} personne du singulier) $\rightarrow$ \eng{goes}. « Chaque dimanche, mon père va pêcher sur le fleuve Nyong. »

2. \eng{do} — \eng{athletics} est une activité individuelle, donc DO ; sujet pluriel \eng{students} $\rightarrow$ \eng{do}. « Les élèves de notre école font de l'athlétisme deux fois par semaine. »

3. \eng{play} — \eng{ludo} est un jeu, donc PLAY ; sujet \eng{My friends and I} (nous) $\rightarrow$ \eng{play}. « Mes amis et moi jouons au ludo pendant la récréation. »

4. \eng{does} — \eng{gymnastics} est une activité individuelle, donc DO ; sujet \eng{Amina} (3\textsuperscript{e} personne du singulier) $\rightarrow$ \eng{does}. « Pendant son temps libre, Amina fait de la gymnastique au stade. »

5. \eng{go} — \eng{swimming} est un verbe en -ing, donc GO ; sujet \eng{We} (nous) $\rightarrow$ \eng{go}. « Nous allons souvent nager à la piscine municipale de Douala. »

6. \eng{play} — pour les instruments de musique, on utilise PLAY : \eng{play the guitar}. Après le modal \eng{can}, pas de \emph{s}. « Mon oncle sait très bien jouer de la guitare. »
\end{exoresolu}

\begin{exercice}[À vous de jouer !]
1. Traduisez en anglais : 
   a) « J'adore regarder des matchs de football à la télévision. » 
   b) « Elle passe son temps libre à écouter de la musique. » 
   c) « Nous sommes passionnés de randonnée. »

2. Trouvez l'erreur dans chaque phrase et corrigez-la : 
   a) \eng{I look TV every evening.} 
   b) \eng{She listens music in her room.} 
   c) \eng{I am keen on to dance.} 
   d) \eng{I bought a novel at the library.}

3. Formez le mot de la même famille (dérivation) : 
   a) recreation (adjectif) ; 
   b) enjoy (nom) ; 
   c) relax (adjectif en -ing).
\end{exercice}

\begin{corrige}[Exercice « À vous de jouer ! »]
1. a) \eng{I am keen on watching football matches on TV.} (ou \eng{I enjoy watching...}) 
   b) \eng{She spends her free time listening to music.} (attention à \eng{listen \textbf{to}}) 
   c) \eng{We are passionate about hiking.} (ou \eng{We are keen on hiking.})

2. a) \eng{I \textbf{watch} TV every evening.} (\eng{watch}, pas \eng{look}) 
   b) \eng{She listens \textbf{to} music in her room.} (préposition \textbf{to} obligatoire) 
   c) \eng{I am keen on \textbf{dancing}.} (verbe en \textbf{-ing} après la préposition \eng{on}) 
   d) \eng{I bought a novel at the \textbf{bookshop}.} (on achète à la librairie \eng{bookshop}, on emprunte à la bibliothèque \eng{library})

3. a) \eng{recreational} ; b) \eng{enjoyment} ; c) \eng{relaxing}.
\end{corrige}

Vous disposez maintenant de tout le lexique et de toutes les constructions pour parler des loisirs. Dans la dernière section, vous mettrez ces outils en pratique en rédigeant un résumé du texte et en donnant votre opinion personnelle.

\coursec{Guided production: summary and opinion}

Dans la section précédente, \emph{Vocabulary: the world of recreation}, nous avons enrichi notre lexique : \eng{hobby}, \eng{leisure}, \eng{to relax}, \eng{board games}, \eng{outdoor activities}, \eng{to unwind}… Il est temps maintenant d'\textbf{utiliser} ce vocabulaire et le contenu du texte \eng{How We Spend Our Free Time} pour produire notre propre écrit. C'est l'objectif de cette dernière étape : apprendre à \cle{résumer} un texte en quelques phrases, à \cle{donner son opinion} correctement, puis à rédiger un paragraphe personnel sur son activité de loisir préférée. Ce sont exactement les compétences exigées à l'épreuve d'anglais du baccalauréat.

\courssub{Summarising a text: the method}

\begin{definition}[What is a summary?]
Un \cle{résumé} (\eng{summary}) est une reformulation \textbf{courte} et \textbf{personnelle} des idées principales d'un texte. On ne copie jamais les phrases de l'auteur : on les \emph{reformule avec ses propres mots} (\eng{in your own words}). Un bon résumé de niveau Terminale compte généralement \textbf{3 à 5 phrases}.
\end{definition}

\begin{methode}[How to write a good summary — 4 steps]
\begin{enumerate}
\item \textbf{Underline one main idea per paragraph.} Relisez le texte et soulignez \emph{une seule} idée essentielle par paragraphe (souvent la première phrase du paragraphe, la \eng{topic sentence}).
\item \textbf{Rephrase without copying.} Reformulez chaque idée avec vos propres mots : changez le vocabulaire (\eng{youngsters} $\rightarrow$ \eng{young people}), changez la structure (voix active $\rightarrow$ voix passive, ou l'inverse).
\item \textbf{Link your sentences with connectors.} Reliez les idées avec des connecteurs logiques : \eng{first}, \eng{then}, \eng{moreover}, \eng{finally}.
\item \textbf{Check the required length.} Vérifiez la longueur demandée (3 à 5 phrases) et supprimez les détails inutiles : exemples, chiffres, anecdotes.
\end{enumerate}
\end{methode}

\begin{propriete}[Connectors for summarising and giving one's opinion]
\begin{itemize}
\item \textbf{Pour résumer (ordre et conclusion)} : \eng{first} (d'abord), \eng{then} (ensuite), \eng{moreover} / \eng{furthermore} (de plus), \eng{finally} (finalement), \eng{in conclusion} / \eng{to sum up} (en conclusion).
\item \textbf{Pour donner son opinion} : \eng{in my opinion} (à mon avis), \eng{I think that} (je pense que), \eng{I believe (that)} (je crois que), \eng{as far as I am concerned} (en ce qui me concerne), \eng{in my view} (selon moi).
\end{itemize}
Ces connecteurs se placent en général \textbf{en tête de phrase}, suivis d'une virgule : \eng{Moreover, many students prefer outdoor games.}
\end{propriete}

\begin{exemplebox}[Opinion phrases: English and French]
\eng{In my opinion, recreation is essential for students.} « À mon avis, les loisirs sont essentiels pour les élèves. »

\eng{As far as I am concerned, football is the best hobby.} « En ce qui me concerne, le football est le meilleur passe-temps. »

\eng{I believe that reading is more relaxing than watching television.} « Je crois que la lecture est plus reposante que la télévision. »

\eng{To sum up, free time helps young people to relax and to socialise.} « En résumé, le temps libre aide les jeunes à se détendre et à socialiser. »
\end{exemplebox}

\courssub{The -ing form after verbs of preference}

Avant de rédiger, révisons une règle de grammaire indispensable pour parler de ses loisirs.

\begin{propriete}[Verb + -ing after enjoy, like, prefer, spend time, be fond of]
Après les verbes \cle{enjoy} (aimer), \cle{like} (aimer), \cle{love} (adorer), \cle{hate} (détester), \cle{prefer} (préférer), \cle{spend time} (passer du temps) et après l'expression \cle{be fond of} (aimer), le verbe qui suit est toujours à la forme en \textbf{-ing} (gérondif) :

\eng{I enjoy reading.} « J'aime lire. »

\eng{She spends her free time dancing.} « Elle passe son temps libre à danser. »

\eng{They are fond of playing board games.} « Ils aiment jouer aux jeux de société. »

\textbf{Remarque} : après \eng{like}, la forme \eng{to + verbe} existe aussi (\eng{I like to read}), mais la forme en \eng{-ing} est la plus sûre et la plus fréquente pour parler de ses goûts habituels.
\end{propriete}

\begin{attention}[Erreurs fréquentes des francophones]
\begin{itemize}
\item \eng{I am agree.} est \textbf{faux} : \eng{agree} est un verbe, pas un adjectif. Dites \eng{I agree.} « Je suis d'accord. » (En français on utilise « être », en anglais non !)
\item \eng{I like play football.} est \textbf{faux} : dites \eng{I like playing football.}
\item N'oubliez pas le \textbf{-s} de la troisième personne du singulier au présent simple : \eng{My brother enjoy\textbf{s} swimming.} « Mon frère aime nager. »
\item Ne \textbf{mélangez pas les temps} : si vous décrivez vos habitudes, restez au présent simple du début à la fin.
\item Ne \textbf{copiez jamais} des phrases entières du texte support dans votre résumé : l'examinateur le sanctionne. Reformulez !
\end{itemize}
\end{attention}

\courssub{Model: a commented summary and opinion}

\begin{textebox}[Model summary and opinion — based on “How We Spend Our Free Time”]
First, the text explains that recreational activities are an important part of family and social life, because they help people to relax after work or school. Then, the writer shows that young people and adults do not spend their free time in the same way: teenagers often prefer sports, music and social media, while older people enjoy quieter activities such as reading or gardening. Moreover, the text points out that some hobbies, like playing football or singing in a choir, bring people together and create friendship. Finally, the writer concludes that free time is precious and that everyone should choose an activity that gives real pleasure.

In my opinion, recreation is essential for students because it reduces stress before examinations. As far as I am concerned, playing basketball with my friends is the best way to spend a Saturday afternoon.
\end{textebox}

\textbf{Commentaire du modèle.} Observez la construction :
\begin{itemize}
\item \textbf{Phrase 1} (\eng{First}) : idée principale du paragraphe 1 — l'importance des loisirs.
\item \textbf{Phrase 2} (\eng{Then}) : idée du paragraphe 2 — la différence jeunes/adultes, reformulée avec \eng{do not spend... in the same way}.
\item \textbf{Phrase 3} (\eng{Moreover}) : idée du paragraphe 3 — les loisirs qui unissent.
\item \textbf{Phrase 4} (\eng{Finally}) : conclusion de l'auteur, reformulée.
\item \textbf{Opinion} : deux phrases personnelles introduites par \eng{In my opinion} et \eng{As far as I am concerned}, avec le verbe en \eng{-ing} après \eng{enjoy} implicite : \eng{playing basketball}.
\end{itemize}

\begin{center}
\begin{tabularx}{\linewidth}{|X|X|X|}
\hline
\rowcolor{popblueL} \textbf{Français} & \textbf{English} & \textbf{Remarque} \\
\hline
D'abord / Premièrement & First / Firstly & En tête de phrase + virgule \\
\hline
Ensuite & Then / Next & \eng{Next} est plus formel \\
\hline
De plus & Moreover / Furthermore & Registre soutenu, idéal à l'écrit \\
\hline
Finalement & Finally / Lastly & Pour la dernière idée \\
\hline
En conclusion & In conclusion / To sum up & Introduit le bilan \\
\hline
À mon avis & In my opinion / In my view & Jamais \eng{in my mind} ! \\
\hline
En ce qui me concerne & As far as I am concerned & Expression très appréciée à l'examen \\
\hline
Je suis d'accord & I agree & Pas de \eng{am} ! \\
\hline
\end{tabularx}
\end{center}

\courssub{Guided production: “My favourite recreational activity”}

Vous allez maintenant rédiger un paragraphe de \textbf{6 à 8 lignes} intitulé \eng{My favourite recreational activity}. Suivez le plan en quatre étapes ci-dessous : il correspond exactement à la structure attendue d'un petit texte d'expression personnelle.

\begin{popfigure}
\begin{tikzpicture}[node distance=0.55cm and 0.9cm,
box/.style={rectangle, rounded corners=3pt, draw=popblue, fill=popblueL, align=center, text width=3.4cm, minimum height=1.1cm, font=\small},
arr/.style={-{Stealth[length=3mm]}, thick, poporange}]
\node[box, align=center] (intro){\textbf{Introduction}\\ Name your favourite activity\\ \emph{My favourite activity is...}};
\node[box, right=of intro, align=center] (desc){\textbf{Description}\\ What? Where? With whom?\\ \emph{I play... with... at...}};
\node[box, right=of desc, align=center] (reasons){\textbf{Reasons}\\ Two arguments\\ \emph{First... Moreover...}};
\node[box, right=of reasons, align=center] (concl){\textbf{Conclusion}\\ Your opinion\\ \emph{In my opinion...}};
\draw[arr] (intro) -- (desc);
\draw[arr] (desc) -- (reasons);
\draw[arr] (reasons) -- (concl);
\end{tikzpicture}
\legende{Structure du paragraphe : Introduction $\rightarrow$ Description $\rightarrow$ Raisons $\rightarrow$ Conclusion.}
\end{popfigure}

\begin{methode}[Writing your paragraph step by step]
\begin{enumerate}
\item \textbf{Introduction (1 phrase)} : nommez votre activité préférée. \eng{My favourite recreational activity is playing handball.}
\item \textbf{Description (2 phrases)} : précisez \emph{quoi}, \emph{où}, \emph{avec qui}, \emph{quand}. \eng{I play every weekend with my classmates on the school playground.}
\item \textbf{Raisons (2 phrases)} : donnez deux arguments avec \eng{first} et \eng{moreover}. \eng{First, it keeps me fit. Moreover, I enjoy spending time with my friends.}
\item \textbf{Conclusion (1 phrase)} : donnez votre opinion. \eng{In my opinion, handball is the best hobby for teenagers.}
\item \textbf{Relecture} : vérifiez le \eng{-s} à la 3\textsuperscript{e} personne, les verbes en \eng{-ing} après \eng{like/enjoy}, et l'unité des temps (présent simple).
\end{enumerate}
\end{methode}

\begin{exoresolu}[Guided writing — solved example]
\textbf{Énoncé.} Rédigez un paragraphe de 6 à 8 lignes intitulé \eng{My favourite recreational activity} en suivant le plan : introduction, description, deux raisons, conclusion avec opinion.

\tcblower
\textbf{Solution détaillée.}

\eng{My favourite recreational activity is reading novels. I usually read in the evening, in my room, after finishing my homework. I also enjoy borrowing books from the school library and discussing them with my best friend. First, reading helps me to relax and to forget the stress of examinations. Moreover, novels teach me new English words and show me how people live in other countries. In my opinion, reading is the most useful hobby for a student because it gives both pleasure and knowledge.}

\textbf{Points de méthode à retenir :}
\begin{itemize}
\item Introduction claire dès la première phrase (\eng{My favourite... is...}).
\item Description complète : quoi (\eng{novels}), quand (\eng{in the evening}), où (\eng{in my room}), avec qui (\eng{my best friend}).
\item Deux raisons introduites par \eng{First} et \eng{Moreover}.
\item Verbes en \eng{-ing} après \eng{enjoy} : \eng{I enjoy borrowing... and discussing...}.
\item Présent simple partout, avec le \eng{-s} : \eng{reading help\textbf{s}}, \eng{novels teach}, \eng{it give\textbf{s}}.
\item Opinion finale introduite par \eng{In my opinion}.
\end{itemize}
\end{exoresolu}

\begin{exercice}[Your turn!]
\begin{enumerate}
\item Résumez le texte \eng{How We Spend Our Free Time} en \textbf{4 phrases} avec vos propres mots, en utilisant \eng{first}, \eng{then}, \eng{moreover} et \eng{in conclusion}.
\item Corrigez les erreurs dans ces phrases :
\begin{enumerate}
\item \eng{I am agree with the writer.}
\item \eng{She enjoy to dance with her friends.}
\item \eng{My father like watching football on television.}
\end{enumerate}
\item Rédigez votre propre paragraphe \eng{My favourite recreational activity} (6 à 8 lignes) en suivant le plan en quatre étapes. Terminez par votre opinion avec \eng{as far as I am concerned}.
\end{enumerate}
\end{exercice}

\begin{corrige}[Exercice — correction des phrases]
\begin{itemize}
\item (a) \eng{I agree with the writer.} — \eng{agree} est un verbe : pas de \eng{am}.
\item (b) \eng{She enjoys dancing with her friends.} — \eng{-s} à la 3\textsuperscript{e} personne et verbe en \eng{-ing} après \eng{enjoy}.
\item (c) \eng{My father likes watching football on television.} — \eng{-s} obligatoire : \eng{likes}.
\end{itemize}
Pour les productions 1 et 3, l'enseignant vérifiera : la reformulation sans copie, la présence des connecteurs, les verbes en \eng{-ing}, le \eng{-s} de la 3\textsuperscript{e} personne et l'unité des temps.
\end{corrige}

\begin{aretenir}[Key points to remember]
\begin{itemize}
\item Un résumé = \textbf{3 à 5 phrases}, une idée par paragraphe, reformulée \emph{avec vos propres mots}.
\item Connecteurs du résumé : \eng{first}, \eng{then}, \eng{moreover}, \eng{finally}, \eng{in conclusion}.
\item Connecteurs de l'opinion : \eng{in my opinion}, \eng{I think that}, \eng{I believe}, \eng{as far as I am concerned}.
\item Après \eng{enjoy}, \eng{like}, \eng{love}, \eng{hate}, \eng{prefer}, \eng{spend time}, \eng{be fond of} : verbe en \textbf{-ing}.
\item Jamais \eng{I am agree} ; toujours \eng{I agree}. Jamais \eng{I like play} ; toujours \eng{I like playing}.
\item Structure d'un paragraphe d'opinion : Introduction $\rightarrow$ Description $\rightarrow$ Raisons $\rightarrow$ Conclusion.
\end{itemize}
\end{aretenir}

\newpage

\coursec{Activité d'intégration}

\coursec{Lesson 7 — Reading: Texts on recreational activities}

\courssub{Objectifs de l'activité}

À la fin de cette activité d'intégration, l'élève sera capable de :
\begin{itemize}
\item mobiliser le \cle{lexique des loisirs} (recreational activities) dans une tâche de communication authentique ;
\item rédiger des \cle{questions} correctes en anglais (ordre des mots, auxiliaires \eng{do/does/did}) ;
\item appliquer les stratégies de lecture étudiées dans la leçon : \cle{skimming} (lire vite pour l'idée générale), \cle{scanning} (repérer une information précise) et \cle{guessing} (déduire le sens des mots inconnus grâce au contexte) ;
\item résumer des données et un texte en quelques phrases ;
\item exprimer une \cle{opinion personnelle} avec des connecteurs adaptés (\eng{in my opinion}, \eng{I think that}, \eng{however}...) ;
\item produire un court \cle{article de journal scolaire} en anglais, seul puis en groupe.
\end{itemize}

\courssub{Situation}

\textbf{Contexte.}\par
Tu es élève en Terminale A au lycée bilingue de Bafoussam. Le club d'anglais (\eng{English Language Club}) prépare le prochain numéro du journal du lycée, \emph{The Lion's Voice}. Le thème du mois est : \eng{“How do young people relax?”}. Ton groupe a été choisi pour réaliser une petite enquête (\eng{survey}) auprès des camarades de la classe et publier un article comparatif.

\textbf{Tâche globale.}\par
Par groupes de quatre, vous allez : (1) rédiger un mini-questionnaire en anglais sur les loisirs préférés de vos camarades ; (2) interroger un autre groupe et noter ses réponses ; (3) lire rapidement le texte ci-dessous sur les loisirs des jeunes Ghanéens ; (4) rédiger ensemble un court article pour \emph{The Lion's Voice} intitulé \eng{“Recreation in our school: survey results”}, qui résume vos résultats, les compare au texte lu et donne votre opinion. Les deux meilleurs articles seront présentés oralement à la classe et publiés dans le journal.

\begin{textebox}[Reading document: “Free time in Accra: what teenagers do after school”]
In Ghana, as in many West African countries, teenagers have less free time than adults often imagine. After long hours at school, many students help their parents at home or in the family shop before they can relax. Yet young Ghanaians still manage to enjoy a rich social life in their free time.

Football is by far the most popular activity. In the evenings, dusty pitches in Accra, Kumasi and Tamale fill with boys — and increasingly girls — playing informal matches. Those who do not play often watch: following the English Premier League on television is almost a national hobby, and arguments about Arsenal, Chelsea and Manchester United can last for hours.

Music and dancing come second. Many teenagers listen to Afrobeats and highlife on their phones, and dancing at birthday parties is a favourite weekend pastime. Social media also takes up a growing share of free time: chatting on WhatsApp and watching short videos are now common ways of unwinding after homework.

Surprisingly, reading for pleasure is less popular than it was twenty years ago, although students who attend reading clubs say it helps them succeed at school. Finally, church activities — choirs, youth groups and drama — occupy many young people on Sundays.

One thing is clear: whatever they do, Ghanaian teenagers prefer relaxing together rather than alone. Free time is, above all, a social affair.
\end{textebox}

\textbf{Glossaire :} \emph{pitch} (n.) : terrain de sport ; \emph{pastime} (n.) : passe-temps ; \emph{to unwind} (v.) : se détendre ; \emph{an affair} (n.) : une affaire, une question.

\courssub{Consignes}

\begin{enumerate}
\item \textbf{Préparer le questionnaire (10 min).} En groupe, rédigez \textbf{cinq questions} en anglais sur les loisirs de vos camarades. Utilisez au moins trois types de questions : une question fermée (\eng{Do you...?}), une question en \emph{wh-} (\eng{What / How often / Where...?}) et une question d'opinion (\eng{Do you think that...?}). Faites vérifier la grammaire par votre professeur ou par un autre groupe.

\item \textbf{Mener l'enquête (10 min).} Interrogez les membres d'un autre groupe. Posez les questions à l'oral, \textbf{en anglais uniquement}, et notez les réponses en quelques mots. Répondez aussi à leurs questions.

\item \textbf{Lire le texte (10 min).} Lisez le document \eng{“Free time in Accra”} en deux étapes :
\begin{itemize}
\item \textbf{Skimming} : lisez le texte en deux minutes et choisissez le meilleur titre alternatif parmi : (a) \eng{“Ghanaian teenagers have no free time”} ; (b) \eng{“How young Ghanaians relax”} ; (c) \eng{“Football in Ghana”}.
\item \textbf{Scanning} : repérez rapidement les réponses à ces questions : \eng{Which sport is the most popular?} — \eng{What do teenagers do on Sundays?} — \eng{Is reading for pleasure more or less popular than before?}
\item \textbf{Guessing} : sans dictionnaire, déduisez le sens de \eng{pastime} et \eng{to unwind} grâce au contexte.
\end{itemize}

\item \textbf{Rédiger l'article (20 min).} Rédigez ensemble un article de \textbf{120 à 150 mots} intitulé \eng{“Recreation in our school: survey results”}. Il doit contenir : une \textbf{introduction} (qui, quoi, où), un \textbf{résumé des résultats} de votre enquête, une \textbf{comparaison} avec le texte sur le Ghana (\eng{similarly}, \eng{unlike}, \eng{whereas}...) et une \textbf{conclusion} donnant l'opinion du groupe.

\item \textbf{Présenter à l'oral (5 min par groupe).} Les deux groupes sélectionnés lisent leur article à la classe. Les autres élèves écoutent et notent une information qu'ils ont apprise (\eng{I learnt that...}).
\end{enumerate}

\courssub{Ressources à mobiliser}

\begin{aretenir}[Boîte à outils linguistique]
\textbf{Lexique des loisirs :} \eng{hobby, pastime, leisure, free time, to relax, to unwind, to hang out with friends, to play football, to watch TV, to listen to music, to chat online, to read for pleasure, dancing, singing, video games}.

\textbf{Former des questions :} auxiliaire + sujet + verbe : \eng{How often do you play football?} ; \eng{Do you prefer reading or watching films?} Attention : pas d'auxiliaire \eng{do} quand le mot interrogatif est sujet : \eng{Who plays football in your family?}

\textbf{Fréquences :} \eng{always, usually, often, sometimes, rarely, never} + \eng{once/twice a week, every day}.

\textbf{Comparaison :} \eng{like, unlike, similarly, whereas, while, both... and..., compared to}.

\textbf{Opinion :} \eng{In my opinion..., I think / I believe that..., It seems to me that..., Personally, I...}.

\textbf{Résumé de résultats :} \eng{Most students..., Half of the group..., Only a few..., The majority of...}.
\end{aretenir}

\begin{attention}[Erreurs fréquentes des francophones]
\begin{itemize}
\item Ne dites pas \eng{*I make sport} mais \eng{I play sport / I do sport} ; pas \eng{*I make my devoirs} mais \eng{I do my homework}.
\item \eng{Actually} ne signifie pas « actuellement » : dites \eng{nowadays} ou \eng{at present}.
\item Ordre des mots dans les questions : \eng{*What you do after school?} est faux ; dites \eng{What do you do after school?}
\item \eng{To assist} ne veut pas dire « assister à » : dites \eng{to attend} (a match, church).
\item Pas de « s » à la 3e personne au négatif : \eng{*He doesn't plays} est faux ; dites \eng{He doesn't play}.
\end{itemize}
\end{attention}

\courssub{Proposition de production et corrigé commenté}

\begin{exoresolu}[Modèle d'article : “Recreation in our school: survey results”]
Rédiger un article de 120 à 150 mots résumant l'enquête de classe, le comparant au texte sur le Ghana et donnant une opinion.
\tcblower
\textbf{Production modèle :}

\eng{Recreation in our school: survey results}

\eng{Last week, our English group interviewed eight classmates about their favourite free-time activities. The results are clear: football and music are the winners.}

\eng{Most students play or watch football at least twice a week, and half of the group listen to music every evening. Chatting with friends on WhatsApp is also very popular. Surprisingly, only a few students read for pleasure, and nobody mentioned dancing.}

\eng{Compared to Ghanaian teenagers, we are very similar: like them, we love football and Afrobeats, and we prefer relaxing with friends rather than alone. However, unlike young Ghanaians, few of us take part in church activities, probably because our weekends are shorter.}

\eng{In our opinion, recreation is not a waste of time: it helps us forget stress and do better at school. We believe the school should create more clubs so that everyone can find an activity they enjoy.}

\textbf{Commentaire des choix de langue :}
\begin{itemize}
\item \textbf{Structure journalistique} : titre, introduction (\eng{Last week, our English group interviewed...} : qui, quoi, quand), résultats, comparaison, opinion. Le lecteur sait immédiatement de quoi parle l'article.
\item \textbf{Résumé des données} sans chiffres inventés compliqués : \eng{most students}, \eng{half of the group}, \eng{only a few}, \eng{nobody} — quantifieurs simples et corrects.
\item \textbf{Connecteurs de comparaison} bien placés : \eng{compared to}, \eng{like them}, \eng{however}, \eng{unlike} + nom. Notez la structure parallèle : \eng{we are very similar: like them, we love...}.
\item \textbf{Opinion explicite et justifiée} : \eng{In our opinion...} suivi d'une justification (\eng{it helps us forget stress}) et d'une suggestion (\eng{the school should create more clubs}) — le modal \eng{should} est idéal pour proposer.
\item \textbf{Grammaire} : présent simple pour les habitudes (\eng{play, listen, is}), prétérit uniquement pour situer l'enquête (\eng{interviewed, mentioned}). Aucune faute d'auxiliaire dans les questions rapportées.
\end{itemize}
\end{exoresolu}

\courssub{Bilan de l'activité}

Cette activité a permis de réinvestir l'ensemble des compétences de la leçon dans une tâche authentique et motivante :
\begin{itemize}
\item \textbf{Compréhension de l'écrit} : le texte sur le Ghana a été lu avec les trois stratégies étudiées (skimming pour l'idée générale, scanning pour les détails, guessing pour le vocabulaire inconnu) ;
\item \textbf{Lexique} : le vocabulaire des loisirs a été utilisé à l'oral (questions/réponses) puis à l'écrit (article) ;
\item \textbf{Grammaire en contexte} : formation des questions, présent simple des habitudes, adverbes de fréquence, connecteurs de comparaison et d'opinion ;
\item \textbf{Expression orale} : conduite d'un entretien en anglais et présentation de l'article devant la classe ;
\item \textbf{Expression écrite} : production d'un article structuré, prêt à être publié dans le journal du lycée.
\end{itemize}

\begin{methode}[Pour aller plus loin : auto-évaluation]
Avant de rendre l'article, chaque groupe se pose quatre questions : (1) Avons-nous un titre et une introduction claire ? (2) Avons-nous résumé nos résultats avec des quantifieurs corrects ? (3) Avons-nous comparé avec le texte ghanéen à l'aide d'au moins deux connecteurs ? (4) Avons-nous donné une opinion justifiée ? Si une réponse est « non », l'article doit être retravaillé avant publication.
\end{methode}

\newpage

\coursec{Méthodes et exercices résolus}

\courssub{Savoir-faire 1 : Lire et comprendre un texte sur les loisirs (skimming, scanning, déduction du sens)}

\begin{methode}[Réussir la compréhension d'un texte sur les loisirs]
\begin{enumerate}
\item \textbf{Lire d'abord le titre et la source.} Le titre (\eng{How We Spend Our Free Time}) annonce le thème : les \cle{recreational activities} (activités de loisirs). Activez votre vocabulaire avant de lire : \eng{hobby, pastime, leisure, entertainment, relaxation}.
\item \textbf{Skimming} (survol) : lisez rapidement le premier et le dernier paragraphe, puis la première phrase de chaque paragraphe, pour trouver l'\cle{idée générale} (\eng{gist}). Ne vous arrêtez pas sur les mots inconnus.
\item \textbf{Lire les questions} avant la deuxième lecture : elles vous disent quoi chercher.
\item \textbf{Scanning} (repérage) : cherchez dans le texte les mots-repères des questions (noms propres, chiffres, verbes d'action) et soulignez la phrase qui contient la réponse.
\item \textbf{Déduire le sens des mots inconnus} (\eng{guessing from context}) : utilisez le contexte (mots voisins, exemples introduits par \eng{such as}, opposition avec \eng{but, unlike, however}), la famille du mot (\eng{relax} $\rightarrow$ \eng{relaxation}) et les préfixes/suffixes (\eng{enjoy-ment}, \eng{use-less}).
\item \textbf{Rédiger la réponse} en phrases complètes, avec vos propres mots, sans recopier tout le paragraphe. Commencez par \eng{According to the text...} ou \eng{The writer says that...}.
\end{enumerate}
\end{methode}

\begin{exoresolu}[Compréhension d'un texte sur les loisirs]
\textbf{Read the text and answer the questions.}

\emph{“In Yaoundé, young people spend their free time in many different ways. Some prefer active pastimes such as playing football in the neighbourhood or jogging at the stadium. Others choose quieter hobbies: reading novels, watching films or chatting with friends online. Whatever the activity, leisure time helps students relax, make friends and develop new skills. However, doctors warn that too much screen time can be harmful to health.”}

\textbf{Questions:}
\begin{enumerate}
\item What is the text about? (skimming)
\item Mention two active pastimes and two quiet hobbies from the text. (scanning)
\item What does the word “pastimes” mean? Guess from the context.
\item According to the text, why is leisure time important?
\end{enumerate}
\tcblower
\textbf{Solution détaillée.}
\begin{enumerate}
\item \textbf{Idée générale (skimming).} La première phrase annonce le sujet : les façons dont les jeunes de Yaoundé occupent leur temps libre. Réponse rédigée : \eng{The text is about how young people in Yaoundé spend their free time.} (« Le texte parle de la façon dont les jeunes de Yaoundé occupent leur temps libre. »)
\item \textbf{Repérage (scanning).} On cherche les mots de la famille des loisirs. Actifs : \eng{playing football}, \eng{jogging}. Calmes : \eng{reading novels}, \eng{watching films}, \eng{chatting with friends online}. Réponse : \eng{Two active pastimes are playing football and jogging. Two quiet hobbies are reading novels and watching films.}
\item \textbf{Déduction du sens.} Le mot \eng{pastimes} est suivi de \eng{such as playing football... jogging}, donc il désigne des activités de loisir ; la composition du mot (\eng{pass} + \eng{time}, « passer le temps ») confirme. Réponse : \eng{“Pastimes” means activities that people do for pleasure in their free time; it is similar to “hobbies”.} (« passe-temps ».)
\item \textbf{Interprétation.} Le texte dit : \eng{leisure time helps students relax, make friends and develop new skills}. Réponse en phrase complète : \eng{According to the text, leisure time is important because it helps students relax, make friends and develop new skills.}
\end{enumerate}
\textbf{Conclusion méthodologique :} une bonne réponse de compréhension = une phrase complète + une information juste + aucun mot recopié inutilement. Justifiez toujours par \eng{according to the text} ou \eng{the writer says that}.
\end{exoresolu}

\courssub{Savoir-faire 2 : Répondre à des questions de compréhension et d'interprétation}

\begin{methode}[Traiter les différents types de questions]
\begin{enumerate}
\item \textbf{Identifier le type de question.} \eng{What / Where / When / Who} = question factuelle (réponse dans le texte). \eng{Why / How} = question d'explication (cause ou manière). \eng{Do you think...? / In your opinion...} = question d'interprétation (votre avis, à justifier).
\item \textbf{Questions factuelles} : repérez le mot-clé de la question dans le texte, lisez la phrase entière, reformulez.
\item \textbf{Questions vrai/faux} (\eng{True or False? Justify}) : citez ou paraphrasez la ligne du texte qui prouve votre choix. Sans justification, la réponse ne vaut que la moitié des points.
\item \textbf{Questions d'interprétation} : donnez votre avis avec \eng{I think that... / In my opinion... / I believe that...} et justifiez avec \eng{because}, \eng{since}, \eng{as}.
\item \textbf{Accorder le temps verbal} de la réponse à celui de la question : question au present simple $\rightarrow$ réponse au present simple ; question au past simple $\rightarrow$ réponse au past simple.
\item \textbf{Relire} : une réponse complète contient un sujet, un verbe conjugué et un complément ; jamais un seul mot.
\end{enumerate}
\end{methode}

\begin{exoresolu}[Questions de compréhension et d'interprétation]
\textbf{Read the passage and answer the questions in complete sentences.}

\emph{“Every weekend, hundreds of supporters gather at the Ahmadou Ahidjo Stadium to watch football matches. For many Cameroonians, football is more than a game: it is a passion that unites families and villages. During important matches, shops close early and the streets become empty.”}

\textbf{Questions:}
\begin{enumerate}
\item Where do supporters gather every weekend?
\item True or False? Justify: “Football is only a game for Cameroonians.”
\item What happens in the streets during important matches?
\item In your opinion, why does football unite people? (2 sentences)
\end{enumerate}
\tcblower
\textbf{Solution détaillée.}
\begin{enumerate}
\item Question factuelle avec \eng{where} $\rightarrow$ on cherche le lieu. Le texte dit \eng{at the Ahmadou Ahidjo Stadium}. Réponse complète : \eng{Every weekend, supporters gather at the Ahmadou Ahidjo Stadium.} (Present simple, car la question est au present simple.)
\item Vrai/faux avec justification. Le texte dit \eng{football is more than a game}, donc l'affirmation est fausse. Réponse : \eng{False. The text says that football is more than a game: it is a passion that unites families and villages.} La justification est indispensable pour avoir tous les points.
\item Question factuelle avec \eng{what}. Réponse : \eng{During important matches, the streets become empty and shops close early.}
\item Question d'interprétation : avis personnel + justification. Modèle de réponse : \eng{In my opinion, football unites people because they share the same emotions when they support their team. Moreover, during a match, differences of age, region and language disappear.} (« À mon avis, le football unit les gens parce qu'ils partagent les mêmes émotions... De plus, pendant un match, les différences d'âge, de région et de langue disparaissent. »)
\end{enumerate}
\textbf{Conclusion :} chaque réponse est une phrase complète, au bon temps, et les questions d'opinion sont toujours justifiées par \eng{because} ou \eng{moreover}.
\end{exoresolu}

\courssub{Savoir-faire 3 : Maîtriser le lexique et les expressions du thème des loisirs}

\begin{methode}[Utiliser correctement le vocabulaire des loisirs]
\begin{enumerate}
\item \textbf{Classer le lexique} en familles : activités sportives (\eng{playing football, jogging, swimming, wrestling}), activités calmes (\eng{reading, watching films, listening to music, playing draughts}), activités sociales (\eng{chatting with friends, dancing, attending ceremonies}).
\item \textbf{Mémoriser les collocations} : on dit \eng{play + sport/jeu} (\eng{play football, play the guitar}), \eng{go + verbe en -ing} (\eng{go swimming, go jogging, go shopping}), \eng{do + activité} (\eng{do gymnastics, do karate}).
\item \textbf{Exprimer les goûts} : \eng{enjoy + V-ing} (\eng{I enjoy reading}), \eng{be fond of / be keen on + V-ing}, \eng{prefer + V-ing}, \eng{hate / dislike + V-ing}.
\item \textbf{Attention aux faux amis} : \eng{actually} = « en fait » (pas « actuellement » = \eng{currently}) ; \eng{sensible} = « raisonnable » (« sensible » = \eng{sensitive}) ; \eng{an opportunity} = « une occasion ».
\item \textbf{Réutiliser le lexique dans une phrase personnelle} : un mot n'est acquis que s'il est employé dans votre propre production.
\end{enumerate}
\end{methode}

\begin{exoresolu}[Exercice de vocabulaire : les loisirs]
\textbf{A. Complete the sentences with the correct word from the box:} \emph{hobby — leisure — pastime — screen time — relax}

1. Reading novels is my favourite \dots\ I do it every evening.
2. Doctors advise students to reduce their \dots\ before sleeping.
3. After a long week at school, I need time to \dots.
4. Football is the most popular \dots\ in Cameroon.
5. What do you do in your \dots\ time?

\textbf{B. Choose the correct verb:} \emph{play / go / do}

6. On Saturdays, we \dots\ swimming at the municipal pool.
7. My sister \dots\ karate at the club.
8. They \dots\ draughts under the mango tree.
\tcblower
\textbf{Solution détaillée.}
\begin{enumerate}
\item \eng{Reading novels is my favourite \textbf{hobby}; I do it every evening.} — \eng{hobby} = « passe-temps favori » ; le possessif \eng{my favourite} appelle un nom d'activité.
\item \eng{Doctors advise students to reduce their \textbf{screen time} before sleeping.} — \eng{screen time} = « temps passé devant les écrans » ; le contexte (\eng{doctors, before sleeping}) oriente vers la santé.
\item \eng{After a long week at school, I need time to \textbf{relax}.} — après \eng{need time to}, il faut un verbe à l'infinitif ; seul \eng{relax} est un verbe.
\item \eng{Football is the most popular \textbf{pastime} in Cameroon.} — \eng{pastime} = « passe-temps, divertissement populaire ».
\item \eng{What do you do in your \textbf{leisure} time?} — collocation figée : \eng{leisure time} = « temps libre ».
\item \eng{On Saturdays, we \textbf{go} swimming at the municipal pool.} — règle : \eng{go + V-ing} pour les activités de déplacement (swimming, jogging, shopping).
\item \eng{My sister \textbf{does} karate at the club.} — \eng{do} pour les arts martiaux et activités individuelles ; attention à la 3\textsuperscript{e} personne : \eng{does}.
\item \eng{They \textbf{play} draughts under the mango tree.} — \eng{play} pour les jeux et les sports avec ballon (\eng{play draughts} = « jouer aux dames »).
\end{enumerate}
\textbf{Conclusion :} le trio \eng{play / go / do} se choisit selon le type d'activité : jeux et sports d'équipe $\rightarrow$ \eng{play} ; activités en \eng{-ing} $\rightarrow$ \eng{go} ; sports individuels et arts martiaux $\rightarrow$ \eng{do}.
\end{exoresolu}

\courssub{Savoir-faire 4 : Résumer un texte et donner son opinion en quelques phrases}

\begin{methode}[Rédiger un résumé suivi d'une opinion]
\begin{enumerate}
\item \textbf{Repérer les 3 ou 4 idées principales} du texte (une par paragraphe en général) et noter les mots-clés.
\item \textbf{Reformuler avec ses propres mots} : ne recopiez jamais de phrases entières du texte ; utilisez des synonymes (\eng{young people} $\rightarrow$ \eng{teenagers} ; \eng{free time} $\rightarrow$ \eng{leisure time}).
\item \textbf{Utiliser des connecteurs} pour lier les idées : \eng{first, then, moreover, however, finally, in conclusion}.
\item \textbf{Respecter la limite de mots} demandée (souvent 40 à 60 mots au Bac) : supprimez les exemples et les détails, gardez l'essentiel.
\item \textbf{Ajouter l'opinion personnelle} en 2 ou 3 phrases : \eng{In my opinion... / I think that... / As far as I am concerned...} + justification avec \eng{because / since / that is why}.
\item \textbf{Relire} : verbes conjugués, accords sujet-verbe (\eng{he enjoy\textbf{s}}), présent simple pour les vérités générales.
\end{enumerate}
\end{methode}

\begin{exoresolu}[Résumé et opinion (type Bac)]
\textbf{Read the text, then summarize it in about 50 words and give your opinion in two or three sentences.}

\emph{“In many Cameroonian towns, leisure habits are changing. In the past, children played outside with their neighbours: they ran, skipped and invented games in the streets. Today, many teenagers prefer to stay indoors. They spend hours on their phones, watching videos or chatting on social networks. Parents and teachers are worried because physical activity is decreasing, and some students no longer know traditional games. Specialists recommend a balance: young people should enjoy technology but also practise sports and outdoor activities to stay healthy.”}
\tcblower
\textbf{Solution détaillée.}

\textbf{Étape 1 — Idées principales :} (1) les habitudes de loisirs changent ; (2) avant : jeux dehors ; aujourd'hui : écrans à l'intérieur ; (3) conséquence : moins d'activité physique, jeux traditionnels oubliés ; (4) conseil : un équilibre entre technologie et sport.

\textbf{Étape 2 — Résumé rédigé (environ 50 mots) :}

\eng{Leisure habits in Cameroon are changing. In the past, children used to play outside, but today many teenagers stay indoors with their phones. As a result, physical activity is decreasing and traditional games are disappearing. Specialists therefore advise young people to balance technology with sports and outdoor activities.} (51 mots)

Remarques : \eng{used to play} exprime une habitude passée révolue ; les connecteurs \eng{but}, \eng{as a result}, \eng{therefore} relient les idées ; aucune phrase n'est recopiée du texte.

\textbf{Étape 3 — Opinion personnelle :}

\eng{In my opinion, the specialists are right: technology is useful, but it should not replace real life. I think that schools should organize sports competitions and teach traditional games so that students can enjoy both worlds.} (« À mon avis, les spécialistes ont raison : la technologie est utile, mais elle ne doit pas remplacer la vie réelle... »)

\textbf{Conclusion :} un bon résumé = idées principales + reformulation + connecteurs + limite de mots respectée ; l'opinion = \eng{in my opinion} + justification.
\end{exoresolu}

\courssub{Savoir-faire 5 : Parler de ses habitudes de loisirs (habitudes présentes et passées)}

\begin{methode}[Exprimer les habitudes : present simple, used to, fréquence]
\begin{enumerate}
\item \textbf{Habitude présente} : present simple + adverbe de fréquence (\eng{always, usually, often, sometimes, rarely, never}). Place de l'adverbe : avant le verbe principal (\eng{I often play football}) mais après \eng{be} (\eng{She is always cheerful}).
\item \textbf{Habitude passée révolue} : \eng{used to + infinitif} (\eng{We used to play hide-and-seek in the village} = « nous jouions à cache-cache »). Négation : \eng{did not use to} ; question : \eng{Did you use to...?}
\item \textbf{Durée} : \eng{for + durée} (\eng{for two hours}), \eng{since + point de départ} (\eng{since 2020}), avec le present perfect si l'action continue (\eng{I have played football since childhood}).
\item \textbf{Comparer passé et présent} : \eng{In the past, children used to..., but nowadays they...}
\item \textbf{Fréquence chiffrée} : \eng{once a week, twice a month, three times a year, every weekend}.
\end{enumerate}
\end{methode}

\begin{exoresolu}[Grammaire en contexte : habitudes de loisirs]
\textbf{A. Put the verbs in brackets in the correct tense.}

1. Every Saturday, my friends and I \dots\ (play) football behind the school.
2. When she was younger, Amina \dots\ (use to / skip) rope with her sisters.
3. I \dots\ (never / watch) television during examination periods.
4. We \dots\ (know) this traditional game since primary school.

\textbf{B. Rewrite the sentence using “used to”.}

5. In the past, children played outside every evening.
\tcblower
\textbf{Solution détaillée.}
\begin{enumerate}
\item \eng{Every Saturday, my friends and I \textbf{play} football behind the school.} — \eng{every Saturday} indique une habitude présente $\rightarrow$ present simple ; sujet pluriel (\eng{my friends and I} = we), donc pas de \eng{-s}.
\item \eng{When she was younger, Amina \textbf{used to skip} rope with her sisters.} — habitude passée révolue (\eng{when she was younger}) $\rightarrow$ \eng{used to + infinitif}. Attention : à la forme affirmative on écrit \eng{used to} ; le \eng{-d} ne disparaît qu'après \eng{did} (\eng{did not use to}).
\item \eng{I \textbf{never watch} television during examination periods.} — l'adverbe de fréquence se place avant le verbe principal ; present simple car habitude.
\item \eng{We \textbf{have known} this traditional game since primary school.} — \eng{since} + point de départ avec une action qui continue $\rightarrow$ present perfect : \eng{have} + participe passé de \eng{know} (\eng{know -- knew -- known}).
\item \eng{In the past, children \textbf{used to play} outside every evening.} — transformation : habitude passée révolue $\rightarrow$ \eng{used to} + base verbale, sans \eng{-ed} au verbe principal.
\end{enumerate}
\textbf{Conclusion :} repérez l'indice temporel (\eng{every Saturday}, \eng{when she was younger}, \eng{since}) : c'est lui qui commande le choix du temps.
\end{exoresolu}

\begin{attention}[Les 5 erreurs qui coûtent des points au Bac]
\begin{enumerate}
\item \textbf{Calque du français pour les goûts.} Ne dites jamais \eng{*I enjoy to play} : après \eng{enjoy}, toujours le \textbf{gérondif} : \eng{I enjoy play\textbf{ing} football.} De même, l'adverbe ne se place pas comme en français : on dit \eng{I like football very much} ou \eng{I very much enjoy reading}, jamais \eng{*I like very much football}.
\item \textbf{Confusion play / go / do.} On dit \eng{play football} (sport avec ballon), \eng{go swimming} (activité en \eng{-ing}), \eng{do karate} (art martial). \eng{*Make sport} est un calque du français « faire du sport » : dites \eng{do sport} ou \eng{play sports}.
\item \textbf{Faux amis.} \eng{Actually} signifie « en fait », pas « actuellement » (\eng{currently}) ; \eng{to assist} ne veut pas dire « assister à » : dites \eng{attend a match}. Écrire \eng{*I assist to the match} coûte deux points (faux ami + préposition en trop).
\item \textbf{Oubli du -s à la 3\textsuperscript{e} personne.} Au present simple : \eng{She enjoy\textbf{s} reading}, \eng{He watch\textbf{es} films}. C'est la faute la plus fréquente dans les résumés et les opinions.
\item \textbf{Orthographe et réponses incomplètes.} Attention à \eng{leisure} (« lé-jour »), \eng{hobby} (un seul \eng{b}), \eng{which} (pas \eng{*wich}). Et répondez toujours aux questions de compréhension par une \textbf{phrase complète} (sujet + verbe conjugué), jamais par un mot isolé : une réponse comme \eng{*At the stadium} ne vaut que la moitié des points.
\end{enumerate}
\end{attention}

\newpage

\coursec{Exercices}

\courssub{Je vérifie mes connaissances}

\begin{exercice}[QCM : stratégies de lecture]
Choisis la bonne réponse pour chaque question.
\begin{enumerate}
\item \emph{Skimming} a text means :
\begin{enumerate}
\item reading every word carefully
\item reading quickly to get the general idea
\item looking for one precise detail
\end{enumerate}
\item \emph{Scanning} a text means :
\begin{enumerate}
\item looking for specific information (a date, a name, a number)
\item translating the whole text
\item reading the text aloud
\end{enumerate}
\item To guess the meaning of an unknown word, you should first :
\begin{enumerate}
\item open your dictionary
\item look at the context and the words around it
\item skip the sentence completely
\end{enumerate}
\item The \emph{topic sentence} of a paragraph is usually :
\begin{enumerate}
\item the last word of the paragraph
\item the sentence that gives the main idea
\item the longest sentence
\end{enumerate}
\end{enumerate}
\end{exercice}

\begin{exercice}[Vrai ou faux, justifié]
Lis le texte suivant, puis dis si les affirmations sont \emph{true} ou \emph{false}. Justifie chaque réponse par une courte citation du texte.
\begin{textebox}[Free time in Yaoundé]
In Yaoundé, young people spend their free time in many different ways. On Saturday afternoons, boys often play football in the neighbourhood, while many girls prefer dancing or singing in church choirs. Video game centres are becoming very popular, but they are expensive, so most teenagers only go there once or twice a month. Watching football matches on television is free in some bars, and it brings people together. Surprisingly, very few young people read books for pleasure, although reading is cheap and relaxing.
\end{textebox}
\begin{enumerate}
\item All teenagers in Yaoundé play football on Saturdays.
\item Video game centres are cheap.
\item Watching football on television can be a social activity.
\item Many young people read books for pleasure.
\end{enumerate}
\end{exercice}

\begin{exercice}[Vocabulaire des loisirs]
Complète chaque phrase avec le mot qui convient, choisi dans la liste : \emph{hobby — leisure — entertaining — spectators — board game — outdoor}.
\begin{enumerate}
\item Gardening is my favourite \ldots\ldots\ldots ; I do it every weekend.
\item Thousands of \ldots\ldots\ldots watched the match at the Ahmadou Ahidjo Stadium.
\item Chess and draughts are my favourite \ldots\ldots\ldots games.
\item Hiking and cycling are \ldots\ldots\ldots activities ; you do them outside.
\item The film was very \ldots\ldots\ldots ; we laughed from the beginning to the end.
\item People have more \ldots\ldots\ldots time during the holidays.
\end{enumerate}
\end{exercice}

\begin{exercice}[Grammaire : play, do ou go ?]
Complète les phrases avec \emph{play}, \emph{do} ou \emph{go}, à la forme qui convient.
\begin{enumerate}
\item Every Sunday, my father \ldots\ldots\ldots draughts with his friends.
\item We usually \ldots\ldots\ldots swimming at the municipal pool in Douala.
\item She \ldots\ldots\ldots judo twice a week ; she is very strong now.
\item Last holidays, they \ldots\ldots\ldots hiking on Mount Cameroon.
\item I \ldots\ldots\ldots not \ldots\ldots\ldots any sport at the moment ; I am too busy.
\end{enumerate}
\end{exercice}

\courssub{Je m'entraîne}

\begin{exercice}[Traduction]
Traduis les phrases suivantes en anglais.
\begin{enumerate}
\item Pendant mon temps libre, j'aime jouer au football avec mes amis.
\item Ma sœur préfère lire des romans parce que c'est relaxant.
\item Beaucoup de jeunes Camerounais passent leurs week-ends à danser.
\item Les jeux de société sont à la fois amusants et éducatifs.
\item Il ne s'intéresse pas du tout aux jeux vidéo ; il les trouve ennuyeux.
\end{enumerate}
\end{exercice}

\begin{exercice}[Transformation de phrases]
Réécris chaque phrase en suivant la consigne entre parenthèses, sans changer le sens.
\begin{enumerate}
\item I like dancing very much. (Use \emph{be fond of}.)
\item She enjoys reading novels. (Turn into a negative sentence.)
\item They play football every Saturday. (Ask a question about the frequency.)
\item Swimming is more relaxing than running. (Use \emph{less \ldots\ than}.)
\item “What do you do in your free time?” the teacher asked me. (Report the question.)
\end{enumerate}
\end{exercice}

\begin{exercice}[Texte à trous]
Complète le texte suivant avec les mots de la liste : \emph{although — take part — relaxing — keen on — spend — such as}.
\begin{textebox}[My weekend]
At the weekend, I do not go to school, so I \ldots\ldots\ldots most of my time with my family. I am very \ldots\ldots\ldots music, and I play the guitar in a small band with my friends. On Saturdays, we \ldots\ldots\ldots in a football tournament in our quarter. \ldots\ldots\ldots I am often tired on Sunday, I never stay in bed late. In the afternoon, I enjoy quiet activities \ldots\ldots\ldots reading or listening to the radio. For me, Sunday evening is the most \ldots\ldots\ldots moment of the week.
\end{textebox}
\end{exercice}

\begin{exercice}[Appariement : mots et définitions]
Associe chaque mot (1--6) à sa définition (a--f).
\begin{center}
\begin{tabularx}{\linewidth}{|l|X|}
\hline
\rowcolor{popblueL} Mot & Définition \\
\hline
1. pastime & a. a person who watches a show or a match \\
\hline
2. spectator & b. an activity you do regularly for pleasure \\
\hline
3. tournament & c. free time, when you do not work or study \\
\hline
4. leisure & d. a competition with several matches or games \\
\hline
5. entertaining & e. amusing and enjoyable \\
\hline
6. choir & f. a group of people who sing together \\
\hline
\end{tabularx}
\end{center}
\end{exercice}

\courssub{Je me prépare au Bac}

\begin{exercice}[Compréhension de texte — type Bac]
Lis attentivement le texte suivant, puis réponds aux questions en anglais, avec tes propres mots autant que possible.
\begin{textebox}[The Afrobeat Festival in Accra]
Every December, thousands of young people travel to Accra, the capital of Ghana, for one of West Africa's biggest music festivals. For three days, the beaches and streets of the city vibrate with Afrobeat, highlife and hip-hop. Famous artists from Nigeria, Ghana and even Cameroon perform on huge open-air stages, and entrance to many concerts is completely free.

The festival is not only about music. During the day, visitors can attend dance workshops, taste Ghanaian dishes such as jollof rice, and buy clothes made by local designers. In the evening, the whole city becomes a giant party.

For many young Ghanaians, the festival is also an economic opportunity. Taxi drivers, food sellers and small hotels earn more money in one week than in two normal months. “December is our best season,” says Ama, who sells grilled fish near the main stage. “Without the festival, my business would be much smaller.”

However, some inhabitants complain about the noise and the mountains of plastic waste left on the beaches. The organisers now promise a “greener” festival, with recycling bins and reusable cups. Whether they will keep this promise remains to be seen.
\end{textebox}
\textbf{Comprehension questions}
\begin{enumerate}
\item Say whether the following statements are \emph{true} or \emph{false}. Justify each answer with a short quotation from the text.
\begin{enumerate}
\item All the concerts at the festival are free.
\item The festival only offers musical activities.
\end{enumerate}
\item Answer the following questions in one or two sentences.
\begin{enumerate}
\item Which countries do the artists come from?
\item Why is December important for Ama's business?
\end{enumerate}
\item Choose the correct answer. Some inhabitants are unhappy because :
\begin{enumerate}
\item the tickets are too expensive
\item the festival creates noise and pollution
\item there are not enough artists
\end{enumerate}
\end{enumerate}
\textbf{Language work}
\begin{enumerate}
\item Find in the text the synonyms of : \emph{well-known} (paragraph 1), \emph{inhabitants} $\to$ \emph{people who live in a place} (paragraph 4), \emph{rubbish} (paragraph 4).
\item Find in the text the opposite of : \emph{to sell} (paragraph 2).
\end{enumerate}
\end{exercice}

\begin{exercice}[Scanning chronométré et déduction lexicale]
\textbf{A. Scanning (3 minutes maximum).} Lis rapidement le texte suivant et retrouve ces six informations : (1) le nom du jeu décrit, (2) le nombre de trous du plateau, (3) le matériau des graines, (4) le moment où l'on y joue, (5) la qualité que le jeu développe, (6) la personne qui enseigne le jeu au narrateur.
\begin{textebox}[Songo, a game of wisdom]
In many Cameroonian villages, evenings end with songo, a traditional board game played by young and old. The board has fourteen holes, and the players move small seeds, usually taken from a local tree, from one hole to another. People play after dinner, when the day's work is finished. My grandfather taught me the rules when I was seven. “Songo is not just a game,” he always says. “It teaches you patience, calculation and respect for your opponent.” Today, tournaments are even organised in some schools in Douala and Bafoussam, and teachers say the game improves children's concentration in mathematics.
\end{textebox}
\textbf{B. Guessing meaning from context.} Sans dictionnaire, devine le sens des mots soulignés dans le texte ci-dessus, puis choisis la bonne traduction française.
\begin{enumerate}
\item \emph{board} : a. planche / plateau de jeu \quad b. tableau d'affichage \quad c. pension
\item \emph{seeds} : a. feuilles \quad b. graines \quad c. pierres précieuses
\item \emph{rules} : a. règles \quad b. règlements de compte \quad c. lignes
\item \emph{opponent} : a. partenaire \quad b. spectateur \quad c. adversaire
\item \emph{improves} : a. empêche \quad b. améliore \quad c. réduit
\item \emph{concentration} : a. concentration \quad b. conversation \quad c. compétition
\end{enumerate}
\end{exercice}

\begin{exercice}[Expression écrite — type Bac]
\textbf{Topic :} Your school magazine has asked students to write about their free time.

Write a paragraph of about \textbf{100 words} about your favourite recreational activity. In your paragraph :
\begin{enumerate}
\item name the activity and say when and where you practise it ;
\item say who you do it with ;
\item explain two reasons why you enjoy it ;
\item say what this activity brings to your life (health, friends, relaxation\ldots).
\end{enumerate}
\textbf{Conseils :} utilise au moins deux expressions du cours (\emph{to be fond of}, \emph{to be keen on}, \emph{in my free time}, \emph{not only \ldots\ but also}) ; vérifie l'accord des verbes à la troisième personne du singulier ; évite de traduire mot à mot depuis le français.
\end{exercice}

\newpage

\coursec{Corrigés des exercices}

\begin{corrige}[Exercice 1 — QCM : stratégies de lecture]
\begin{enumerate}
\item \textbf{b.} reading quickly to get the general idea. \emph{Skimming} consiste à lire rapidement (titre, première phrase de chaque paragraphe) pour saisir l'idée générale du texte.
\item \textbf{a.} looking for specific information (a date, a name, a number). \emph{Scanning} consiste à parcourir le texte des yeux pour repérer une information précise, sans tout lire.
\item \textbf{b.} look at the context and the words around it. Le contexte (mots voisins, grammaire de la phrase, logique du paragraphe) permet de déduire le sens avant de recourir au dictionnaire.
\item \textbf{b.} the sentence that gives the main idea. La \emph{topic sentence}, souvent placée au début du paragraphe, annonce l'idée principale.
\end{enumerate}
\end{corrige}

\begin{corrige}[Exercice 2 — Vrai ou faux, justifié]
\begin{enumerate}
\item \textbf{False.} The text says “boys \emph{often} play football”, not all teenagers ; moreover, “many girls prefer dancing or singing”.
\item \textbf{False.} The text says video game centres “are expensive”.
\item \textbf{True.} The text says watching matches on television “brings people together”, which means it is a social activity.
\item \textbf{False.} The text says “very few young people read books for pleasure”.
\end{enumerate}
\textbf{Point de méthode :} attention aux mots comme \emph{all}, \emph{always}, \emph{never} dans les affirmations : le texte dit souvent \emph{often}, \emph{many}, \emph{very few}, ce qui rend l'affirmation fausse.
\end{corrige}

\begin{corrige}[Exercice 3 — Vocabulaire des loisirs]
\begin{enumerate}
\item Gardening is my favourite \textbf{hobby} ; I do it every weekend. (\emph{hobby} = passe-temps pratiqué régulièrement)
\item Thousands of \textbf{spectators} watched the match at the Ahmadou Ahidjo Stadium. (\emph{spectators} = spectateurs, au pluriel avec \emph{thousands of})
\item Chess and draughts are my favourite \textbf{board} games. (\emph{board games} = jeux de société)
\item Hiking and cycling are \textbf{outdoor} activities ; you do them outside. (\emph{outdoor} = en plein air)
\item The film was very \textbf{entertaining} ; we laughed from the beginning to the end. (\emph{entertaining} = divertissant)
\item People have more \textbf{leisure} time during the holidays. (\emph{leisure time} = temps libre)
\end{enumerate}
\end{corrige}

\begin{corrige}[Exercice 4 — Grammaire : play, do ou go ?]
\begin{enumerate}
\item Every Sunday, my father \textbf{plays} draughts with his friends. (\emph{play} + jeu/sport avec ballon ou plateau ; \emph{-s} à la 3\textsuperscript{e} personne)
\item We usually \textbf{go} swimming at the municipal pool in Douala. (\emph{go} + verbe en \emph{-ing})
\item She \textbf{does} judo twice a week ; she is very strong now. (\emph{do} + sport individuel/de combat ; \emph{does} à la 3\textsuperscript{e} personne)
\item Last holidays, they \textbf{went} hiking on Mount Cameroon. (prétérit irrégulier de \emph{go} : \emph{went})
\item I \textbf{do} not \textbf{do} any sport at the moment ; I am too busy. (auxiliaire \emph{do} pour la négation + \emph{do} comme verbe principal : \emph{do sport})
\end{enumerate}
\textbf{Rappel :} \emph{play} + sport d'équipe/jeu (play football, play the guitar), \emph{do} + activité individuelle (do judo, do gymnastics), \emph{go} + activité en \emph{-ing} (go swimming, go hiking).
\end{corrige}

\begin{corrige}[Exercice 5 — Traduction]
\begin{enumerate}
\item \eng{In my free time, I like playing football with my friends.} (ou \emph{I enjoy playing\ldots} ; \emph{like} + \emph{-ing})
\item \eng{My sister prefers reading novels because it is relaxing.} (ou \emph{prefers to read} ; \emph{relaxing} = relaxant, et non \emph{relaxed} qui décrit la personne)
\item \eng{Many young Cameroonians spend their weekends dancing.} (structure : \emph{spend} + temps + \emph{-ing})
\item \eng{Board games are both fun and educational.} (ou \emph{entertaining and educational})
\item \eng{He is not interested in video games at all ; he finds them boring.} (\emph{to be interested in} ; \emph{boring} = ennuyeux, \emph{bored} = ennuyé)
\end{enumerate}
\textbf{Points de vigilance :} ne pas écrire « during my free time » (calque) mais \emph{in my free time} ; « passer du temps à faire » se traduit par \emph{spend time doing}, jamais par \emph{pass}.
\end{corrige}

\begin{corrige}[Exercice 6 — Transformation de phrases]
\begin{enumerate}
\item \eng{I am very fond of dancing.} (\emph{be fond of} + \emph{-ing} ; \emph{very much} devient \emph{very} avant l'adjectif)
\item \eng{She does not enjoy reading novels.} (négation au présent simple : \emph{does not} + base verbale)
\item \eng{How often do they play football?} (\emph{How often} interroge sur la fréquence)
\item \eng{Running is less relaxing than swimming.} (on inverse les termes avec \emph{less \ldots\ than})
\item \eng{The teacher asked me what I did in my free time.} (discours indirect : pas d'inversion, \emph{do} $\to$ \emph{did}, suppression du point d'interrogation)
\end{enumerate}
\end{corrige}

\begin{corrige}[Exercice 7 — Texte à trous]
At the weekend, I do not go to school, so I \textbf{spend} most of my time with my family. I am very \textbf{keen on} music, and I play the guitar in a small band with my friends. On Saturdays, we \textbf{take part} in a football tournament in our quarter. \textbf{Although} I am often tired on Sunday, I never stay in bed late. In the afternoon, I enjoy quiet activities \textbf{such as} reading or listening to the radio. For me, Sunday evening is the most \textbf{relaxing} moment of the week.
\begin{itemize}
\item \emph{be keen on} = aimer beaucoup ; \emph{take part in} = participer à ; \emph{although} = bien que (en tête de phrase, suivi d'une virgule) ; \emph{such as} = tel que ; \emph{relaxing} = relaxant.
\end{itemize}
\end{corrige}

\begin{corrige}[Exercice 8 — Appariement : mots et définitions]
\begin{center}
\begin{tabular}{|l|l|}
\hline
\rowcolor{popblueL} Mot & Réponse \\
\hline
1. pastime & b. an activity you do regularly for pleasure \\
\hline
2. spectator & a. a person who watches a show or a match \\
\hline
3. tournament & d. a competition with several matches or games \\
\hline
4. leisure & c. free time, when you do not work or study \\
\hline
5. entertaining & e. amusing and enjoyable \\
\hline
6. choir & f. a group of people who sing together \\
\hline
\end{tabular}
\end{center}
\textbf{Astuce :} \emph{spectator} contient la racine latine \emph{spectare} (regarder), comme \emph{spectacle} ; \emph{choir} se prononce « kouaïeur » et désigne une chorale.
\end{corrige}

\begin{corrige}[Exercice 9 — Compréhension de texte — type Bac]
\textbf{Comprehension questions}
\begin{enumerate}
\item \begin{enumerate}
\item \textbf{False.} The text says “entrance to \emph{many} concerts is completely free”, not all of them.
\item \textbf{False.} The text says “The festival is \emph{not only} about music” and mentions “dance workshops”, food and clothes.
\end{enumerate}
\item \begin{enumerate}
\item \eng{The artists come from Nigeria, Ghana and Cameroon.}
\item \eng{December is her best season because the festival brings many visitors, so she sells much more grilled fish than usual.}
\end{enumerate}
\item \textbf{b.} the festival creates noise and pollution. (The text mentions “the noise and the mountains of plastic waste”.)
\end{enumerate}
\textbf{Language work}
\begin{enumerate}
\item \emph{well-known} = \textbf{famous} (paragraph 1) ; \emph{people who live in a place} = \textbf{inhabitants} (paragraph 4) ; \emph{rubbish} = \textbf{waste} (paragraph 4).
\item The opposite of \emph{to sell} is \textbf{to buy} (paragraph 2 : “buy clothes made by local designers”).
\end{enumerate}
\textbf{Barème indicatif (sur 10) :} Vrai/Faux justifiés : 4 pts (1 pt par réponse + justification) ; questions ouvertes : 3 pts ; QCM : 1 pt ; travail lexical : 2 pts.
\textbf{Points de vigilance :} une justification sans citation exacte ne vaut pas tous les points ; répondre par des phrases complètes ; ne pas recopier de longs passages entiers.
\end{corrige}

\begin{corrige}[Exercice 10 — Scanning chronométré et déduction lexicale]
\textbf{A. Scanning}
\begin{enumerate}
\item The name of the game : \textbf{songo}.
\item The number of holes : \textbf{fourteen}.
\item The material of the seeds : they are \textbf{taken from a local tree}.
\item When people play : \textbf{after dinner / in the evening}, when the day's work is finished.
\item The quality the game develops : \textbf{patience} (also calculation and respect ; it improves children's \textbf{concentration} in mathematics).
\item Who taught the narrator : \textbf{his grandfather}, when he was seven.
\end{enumerate}
\textbf{B. Guessing meaning from context}
\begin{enumerate}
\item \emph{board} : \textbf{a.} plateau de jeu (contexte : “a traditional board game” avec des trous)
\item \emph{seeds} : \textbf{b.} graines (contexte : “taken from a local tree”)
\item \emph{rules} : \textbf{a.} règles (contexte : “taught me the rules” d'un jeu)
\item \emph{opponent} : \textbf{c.} adversaire (contexte : “respect for your opponent” dans une partie)
\item \emph{improves} : \textbf{b.} améliore (contexte positif : les enseignants louent le jeu)
\item \emph{concentration} : \textbf{a.} concentration (faux ami évident avec “conversation” : contexte des mathématiques)
\end{enumerate}
\textbf{Point de méthode :} pour deviner, pose-toi trois questions : quelle est la nature grammaticale du mot ? que dit la phrase autour ? quel sens est logique avec le sujet du texte ?
\end{corrige}

\begin{corrige}[Exercice 11 — Expression écrite — type Bac]
\textbf{Proposition de rédaction modèle (environ 100 mots) :}
\begin{textebox}[Model paragraph]
My favourite recreational activity is playing basketball. In my free time, especially on Wednesday and Saturday afternoons, I go to the school playground with my classmates. I am really fond of this sport for two main reasons. First, it keeps me fit and healthy because we run and jump a lot. Secondly, basketball is a team game, so it teaches me to cooperate with others and to respect my opponents. I am also keen on watching professional matches on television with my father. This activity brings me not only good health but also true friends, and after a game I always feel relaxed and happy.
\end{textebox}
\textbf{Traduction :} « Mon activité de loisir préférée est le basket-ball. Pendant mon temps libre, surtout les mercredis et samedis après-midi, je vais au terrain de l'école avec mes camarades. J'aime vraiment ce sport pour deux raisons principales. D'abord, il me garde en forme et en bonne santé parce que nous courons et sautons beaucoup. Ensuite, le basket est un sport d'équipe, donc il m'apprend à coopérer avec les autres et à respecter mes adversaires. J'aime aussi beaucoup regarder des matchs professionnels à la télévision avec mon père. Cette activité m'apporte non seulement une bonne santé mais aussi de vrais amis, et après un match je me sens toujours détendu et heureux. »
\textbf{Barème indicatif (sur 10) :} respect de la tâche et des 4 consignes : 4 pts ; richesse du vocabulaire des loisirs et expressions du cours : 2 pts ; correction grammaticale (accords, 3\textsuperscript{e} personne, \emph{-ing}) : 3 pts ; organisation et longueur (environ 100 mots) : 1 pt.
\textbf{Points de vigilance :}
\begin{itemize}
\item \emph{play basketball} sans article ; mais \emph{play \textbf{the} guitar} (instruments de musique).
\item Après \emph{enjoy}, \emph{be fond of}, \emph{be keen on} : verbe en \emph{-ing} (\emph{I enjoy playing}, jamais « I enjoy to play »).
\item Ne pas traduire « cela me détend » par « it relaxes me » de façon maladroite : préférer \emph{it helps me relax} ou \emph{I feel relaxed}.
\item Éviter les répétitions de \emph{very much} ; varier avec \emph{really}, \emph{especially}, \emph{not only \ldots\ but also}.
\end{itemize}
\end{corrige}

\newpage

\coursec{Fiche bilan}

\begin{popfigure}
\begin{center}\tcbox[colback=popink,colframe=popink,coltext=white,arc=9pt,boxsep=4pt,left=6pt,right=6pt]{\bfseries\small Reading: Recreational Activities}\end{center}
\vspace{-2pt}
\begin{tcbraster}[raster columns=3,raster equal height=rows,raster column skip=6pt,raster row skip=6pt,enhanced,blankest]
\begin{tcolorbox}[enhanced,colback=white,colframe=poporange,boxrule=0.9pt,arc=3pt,title={Key vocabulary},fonttitle=\bfseries\scriptsize,coltitle=white,colbacktitle=poporange,left=4pt,right=4pt,top=2pt,bottom=2pt,boxsep=1pt,toptitle=1.5pt,bottomtitle=1.5pt,halign=left]
\begin{itemize}[nosep,leftmargin=*,itemsep=1pt,font=\scriptsize]
\item leisure, hobby, pastime
\item indoor / outdoor
\end{itemize}
\end{tcolorbox}
\begin{tcolorbox}[enhanced,colback=white,colframe=popgreen,boxrule=0.9pt,arc=3pt,title={Skimming},fonttitle=\bfseries\scriptsize,coltitle=white,colbacktitle=popgreen,left=4pt,right=4pt,top=2pt,bottom=2pt,boxsep=1pt,toptitle=1.5pt,bottomtitle=1.5pt,halign=left]
\begin{itemize}[nosep,leftmargin=*,itemsep=1pt,font=\scriptsize]
\item title + first lines
\end{itemize}
\end{tcolorbox}
\begin{tcolorbox}[enhanced,colback=white,colframe=poppurple,boxrule=0.9pt,arc=3pt,title={Scanning},fonttitle=\bfseries\scriptsize,coltitle=white,colbacktitle=poppurple,left=4pt,right=4pt,top=2pt,bottom=2pt,boxsep=1pt,toptitle=1.5pt,bottomtitle=1.5pt,halign=left]
\begin{itemize}[nosep,leftmargin=*,itemsep=1pt,font=\scriptsize]
\item names, dates, numbers
\end{itemize}
\end{tcolorbox}
\begin{tcolorbox}[enhanced,colback=white,colframe=poppink,boxrule=0.9pt,arc=3pt,title={Guessing meaning},fonttitle=\bfseries\scriptsize,coltitle=white,colbacktitle=poppink,left=4pt,right=4pt,top=2pt,bottom=2pt,boxsep=1pt,toptitle=1.5pt,bottomtitle=1.5pt,halign=left]
\begin{itemize}[nosep,leftmargin=*,itemsep=1pt,font=\scriptsize]
\item context + word family
\end{itemize}
\end{tcolorbox}
\begin{tcolorbox}[enhanced,colback=white,colframe=popgold,boxrule=0.9pt,arc=3pt,title={Comprehension},fonttitle=\bfseries\scriptsize,coltitle=white,colbacktitle=popgold,left=4pt,right=4pt,top=2pt,bottom=2pt,boxsep=1pt,toptitle=1.5pt,bottomtitle=1.5pt,halign=left]
\begin{itemize}[nosep,leftmargin=*,itemsep=1pt,font=\scriptsize]
\item true/false, wh-questions
\end{itemize}
\end{tcolorbox}
\begin{tcolorbox}[enhanced,colback=white,colframe=popblue!70,boxrule=0.9pt,arc=3pt,title={Summary \& opinion},fonttitle=\bfseries\scriptsize,coltitle=white,colbacktitle=popblue!70,left=4pt,right=4pt,top=2pt,bottom=2pt,boxsep=1pt,toptitle=1.5pt,bottomtitle=1.5pt,halign=left]
\begin{itemize}[nosep,leftmargin=*,itemsep=1pt,font=\scriptsize]
\item in my opinion...
\end{itemize}
\end{tcolorbox}
\end{tcbraster}

\legende{Carte mentale de la leçon : lire un texte sur les loisirs}
\end{popfigure}

\begin{bilan}
\textbf{1. Lexique clé : le monde des loisirs (recreation)}
\begin{center}
\begin{tabularx}{\linewidth}{|X|X|}
\hline
\rowcolor{popblueL} English & Français \\
\hline
leisure / free time / spare time & loisirs / temps libre \\
\hline
hobby, pastime & passe-temps, hobby \\
\hline
recreational activity & activité de loisir \\
\hline
indoor / outdoor activities & activités d'intérieur / de plein air \\
\hline
to relax, to unwind & se détendre \\
\hline
to hang out with friends & traîner / sortir avec des amis \\
\hline
to take up a hobby & se mettre à un hobby \\
\hline
to be fond of / keen on (sport) & être passionné de (sport) \\
\hline
entertainment & divertissement \\
\hline
to watch a match / to play draughts & regarder un match / jouer aux dames \\
\hline
\end{tabularx}
\end{center}

\textbf{2. Stratégies de lecture à connaître par c\oe ur}
\begin{center}
\begin{tabularx}{\linewidth}{|l|X|X|}
\hline
\rowcolor{popblueL} Stratégie & Définition & Quand l'utiliser \\
\hline
\cle{Skimming} & Lire vite pour saisir l'idée générale (titre, 1\iere{} et dernière phrases des paragraphes). & Avant les questions : « What is the text about? » \\
\hline
\cle{Scanning} & Chercher un détail précis (nom, date, chiffre, lieu) sans tout relire. & Questions factuelles : who, when, where, how many. \\
\hline
\cle{Guessing} & Déduire le sens d'un mot inconnu grâce au contexte, à la nature du mot et à sa famille. & Mots difficiles soulignés dans le texte. \\
\hline
\end{tabularx}
\end{center}

\textbf{3. Grammaire utile pour parler des loisirs}
\begin{itemize}
\item \cle{like / love / enjoy / hate + V-ing} : \eng{I enjoy playing football.} « J'aime jouer au football. » (jamais \emph{enjoy to play}).
\item \cle{be fond of / be keen on / be interested in + V-ing} : \eng{She is keen on reading novels.} « Elle adore lire des romans. »
\item Préférence : \eng{I prefer dancing to singing.} « Je préfère danser plutôt que chanter. » ; \eng{I would rather stay at home.} « Je préférerais rester à la maison. »
\item Fréquence : \eng{always, usually, often, sometimes, rarely, never} + présent simple : \eng{We often watch football on Saturdays.}
\end{itemize}

\textbf{4. Méthode express : résumer et donner son opinion}
\begin{enumerate}
\item Identifier le thème : \eng{The text is about how Cameroonians spend their free time.}
\item Relever 2 ou 3 idées principales, sans recopier de phrases entières.
\item Reformuler avec ses propres mots (paraphrase).
\item Donner son avis : \eng{In my opinion... / I think that... / As far as I am concerned...}
\item Conclure en une phrase.
\end{enumerate}
\end{bilan}

\begin{aretenir}[Checklist avant le Bac]
\begin{itemize}
\item[$\square$] Je sais définir et distinguer \eng{skimming} et \eng{scanning}.
\item[$\square$] Je lis d'abord le titre et la première phrase de chaque paragraphe pour trouver l'idée générale.
\item[$\square$] Je sais repérer rapidement un nom, une date ou un chiffre dans le texte.
\item[$\square$] Je devine le sens d'un mot inconnu grâce au contexte avant de chercher dans le glossaire.
\item[$\square$] Je connais le lexique des loisirs : \eng{hobby, leisure, pastime, indoor, outdoor, entertainment}.
\item[$\square$] J'utilise correctement \eng{enjoy / like / be keen on + V-ing} sans mettre d'infinitif.
\item[$\square$] Je justifie mes réponses aux questions de compréhension par un élément du texte.
\item[$\square$] Je réponds aux questions vrai/faux en citant la ligne ou la phrase qui justifie ma réponse.
\item[$\square$] Je sais résumer un texte en 3 à 5 phrases avec mes propres mots.
\item[$\square$] Je sais exprimer mon opinion : \eng{In my opinion, reading is more rewarding than watching TV.}
\item[$\square$] Je relis ma production pour vérifier l'accord sujet-verbe et les prépositions (\eng{fond of, keen on, interested in}).
\item[$\square$] Je gère mon temps : 5 minutes de survol, puis les questions, puis la relecture.
\end{itemize}
\end{aretenir}

\findecours

\end{document}
